\documentclass[12pt]{amsart}
\usepackage{amsthm,amsfonts,amssymb,amscd,mathrsfs}
\usepackage{bezier,longtable,amstext}
\usepackage{amssymb}

\sloppy
\topmargin-10mm
\oddsidemargin-5mm
\evensidemargin-5mm
\setlength{\textwidth}{17cm}
\setlength{\textheight}{25cm}
\usepackage[all]{xy}


\newcommand{\EXT}{{\mathcal Ext}}
\newcommand\uC{{\underline{C}}}
\newcommand\Aut{\text{Aut}}
\newcommand\codim{\text{codim}}
\newcommand\coker{\operatorname{coker}\nolimits}
\newcommand{\const}{\operatorname{const}\nolimits}
\newcommand\ev{\text{ev}}
\newcommand\ext{\text{ext}}
\newcommand\Ext{\operatorname{Ext}\nolimits}
\newcommand\GL{\text{GL}}
\newcommand{\Hilb}{\operatorname{Hilb}\nolimits}
\newcommand\hhom{\text{hom}}
\newcommand\Hom{\operatorname{Hom}\nolimits}
\newcommand\id{\text{id}}
\newcommand\im{\operatorname{Im}\nolimits}
\newcommand\Ker{\operatorname{Ker}\nolimits}
\newcommand\length{\text{length}}
\newcommand\mult{\text{mult}}
\newcommand\odd{\text{odd}}
\newcommand{\Pic}{\operatorname{Pic}\nolimits}
\newcommand\pr{\text{pr}}
\newcommand{\Proj}{\operatorname{Proj}\nolimits}
\newcommand\red{\text{red}}
\newcommand\res{\text{res}}
\newcommand\rk{\text{rk}}
\newcommand\Fl{\text{Fl}}
\newcommand{\Sing}{\operatorname{Sing}\nolimits}
\newcommand\Span{\text{Span}}
\newcommand\Stab{\text{Stab}}
\newcommand\Sym{\operatorname{Sym}\nolimits}
\renewcommand{\phi}{\varphi}

\newcommand{\refle}[1]{Lemma \ref{#1}}
\newcommand{\refth}[1]{Theorem \ref{#1}}
\newcommand{\refcor}[1]{Corollary \ref{#1}}
\newcommand{\refprop}[1]{Proposition \ref{#1}}
\newcommand{\refsec}[1]{Section~ \ref{#1}}
\newcommand{\refeq}[1]{(\ref{#1})}

\newcommand {\eqdef} {:=}
\newcommand {\eqdefrev} {=:}

\newcommand\elm{{\operatorname{elm}\nolimits}}
\newcommand\Supp{{\operatorname{Supp}\nolimits}}

\newtheorem{theorem}{Theorem}[section]
\newtheorem{proposition}[theorem]{Proposition}
\newtheorem{prop}[theorem]{Proposition}
\newtheorem{conjecture}[theorem]{Conjecture}
\newtheorem{lemma}[theorem]{Lemma}
\newtheorem{corollary}[theorem]{Corollary}
\newtheorem{sub}[subsection]{}
\newtheorem{subsub}{}[subsection]
\newtheorem{subsubsub}{}[subsubsection]
\newtheorem{assumption}[theorem]{Assumption}
\newtheorem{claim}[theorem]{Claim}
\newtheorem{notation}[theorem]{Notation}


\theoremstyle{definition}
\newtheorem{definition}[theorem]{Definition}
\newtheorem{example}[theorem]{Example}
\newtheorem{proposition-definition}[theorem]{Proposition-Definition}

\newtheorem{remark}[theorem]{Remark}


\nonstopmode



\begin{document}

\begin{flushright}
To the memory of Andrey Todorov
\end{flushright}

\title{Linear ind-grassmannians}

\author[I.~Penkov]{\;Ivan~Penkov}

\vspace{1cm}

\address{
Jacobs University Bremen \\
School of Engineering and Science, Campus Ring 1, 28759 Bremen, Germany}
\email{i.penkov@jacobs-university.de}

\author[A.~Tikhomirov]{\;Alexander S.~Tikhomirov}


\address{
Department of Mathematics\\
State Pedagogical University\\
Respublikanskaya Str. 108
\newline 150 000 Yaroslavl, Russia}
\email{astikhomirov@mail.ru}







\maketitle


\thispagestyle{empty}

\begin{abstract}
We consider ind-varieties obtained as direct limits of chains of embeddings
$X_1\stackrel{\phi_1}{\hookrightarrow}\dots\stackrel{\phi_{m-1}}{\hookrightarrow}
X_m\stackrel{\phi_m}{\hookrightarrow}X_{m+1}\stackrel{\phi_{m+1}}{\hookrightarrow}\dots$,
where each $X_m$ is a grassmannian or an isotropic grassmannian (possibly mixing grassmannians and isotropic
grassmannians), and the embeddings $\phi_m$ are linear in the sense that they induce
isomorphisms of Picard groups. We prove that any such ind-variety is isomorphic to one of certain standard
ind-grassmannians and that the latter are pairwise non-isomorphic ind-varieties.



2010 Mathematics Subject Classification, Primary 14A10, 14M15.

Bibliography: 8 items.

\keywords{Keywords: grassmannian, ind-variety, linear morphism of algebraic varieties.}

\thanks{}
\end{abstract}

\section{Introduction}\label{sec1}
\label{Introduction}
The Barth--Van de Ven--Tyurin--Sato Theorem claims that any finite rank vector bundle on the infinite complex
projective space $\mathbf{P}^\infty$ is isomorphic to a direct sum  of line bundles. For rank two bundles this theorem
has been proved by Barth and Van de Ven in \cite{BV}, and in the general case the theorem has been proved by Tyurin in
\cite{T} and Sato in [S1]. In the last decade we have studied more general ind-varieties for which the result
holds true \cite{PT1}, \cite{PT2} \cite{DP}.

This study has naturally led us to the problem of constructing non-isomorphic ind-varieties arising as direct limits
of given classes of embeddings of projective varieties. In the present
note we address a classification problem along those lines:
we consider linear embeddings of grassmannians, i.e.
embeddings $i:X_1\hookrightarrow X_2$ of a grassmannian $X_1$ into a grassmannian $X_2$ satisfying the condition
$i^*\mathcal{O}_{X_2}(1)\simeq\mathcal{O}_{X_1}(1)$, and determine
how many non-isomorphic ind-varieties can be obtained from such embeddings. Moreover,
we consider also orthogonal and symplectic grassmannians (i.e. isotropic grassmannians arising from non-degenerate
orthogonal or
symplectic forms) and define a \textit{linear ind-grassmannian} as an ind-variety arising as the direct limit
$\underset{\longrightarrow}\lim\ X_n$ of any chain of linear embeddings
$$
X_1\hookrightarrow X_2\hookrightarrow\dots\hookrightarrow X_m \hookrightarrow X_{m+1}\hookrightarrow\dots
$$
of grassmannians, some or all of them orthogonal or symplectic.

Our main result (Theorem 2, see Section \ref{sec5}) states that each linear ind-grassmannian is isomorphic (as an
ind-variety) to one of the standard ind-grassmannians introduced in \cite{DiP}. In particular, any linear
ind-grassmannian is a homogeneous space
of one of the three classical ind-groups $SL(\infty)$, $O(\infty)$, $Sp(\infty)$. We also prove in Theorem 2 that the
standard ind-grassmannians are pairwise non-isomorphic. To make the note self-contained, we do not rely on the
article \cite{DiP}, but introduce (in Section \ref{stand Gr} below) the standard ind-grassmannians in terms of
explicit chains of embeddings.

The main tool we use in Theorem 2 is Theorem 1 (see Section \ref{Lin emb}) which describes linear morphisms of
grassmannians, as well as isotropic grassmannians.

In the related paper \cite{PT3} we return to the original question of extending the generality of the Barth-Van de Ven-
Tyurin-Sato theorem. There we give the list of linear ind-grassmannians on which a bundle of finite rank is isomorphic
to a direct sum of line bundles.

\textbf{Acknowledgement.} A.S.T. was been financially supported by Ministry of Education and Science
of Russian Federation, and acknowledges partial support from Jacobs University Bremen. Both A.S.T. and I.P.
acknowledge the hospitality of the Max-Planck-Institute for Mathematics in Bonn. In addition, I.P. acknowledges
partial DFG support through "DFG Schwerpunkt 1388, Darstellungstheorie".



\vspace{1cm}

\section{Preliminaries}\label{sec2}
\vspace{0.5cm}

\begin{sub}{\bf Notation and conventions.}\label{subsection 2.1}
\rm
Recall that $\mathbb{N}=\{0,1,2,...\}$. We set $\mathbb{Z}_+=\{1,2,3,...\}$. All vector spaces and algebraic varieties
are defined over an algebraically closed field $\mathbb{F}$ of characteristic 0.
The superscript ${}^*$ indicates dual space or dual vector bundle as well as inverse image. If $X$ is a projective
variety with Picard group isomorphic to $\mathbb{Z}$, then $\mathcal{O}_X(1)$ stands for the ample generator of the
Picard group.

By $G(k,V),\ 1\le k\le\dim V,$ we denote the grassmannian of $k$-dimensional subspaces of a finite-dimensional vector
space $V$. For $k=1$, $G(k,V)=\mathbb{P}(V)$. Furthermore, $\mathcal{O}_{G(k,V)}(1)\cong\wedge^kS_{G(k,V)}^*$, where
$S_{G(k,V)}$ is the tautological bundle on $G(k,V)$,  and $\Pic G(k,V)\cong\mathbb{Z}\mathcal{O}_{G(k,V)}(1)$.


In what follows we will consider, both symmetric and symplectic, quadratic forms $\Phi$ on $V$.
Under the assumption that $\Phi$ is fixed, we set $W^\bot:=\{v\in V|\ \Phi(v,w)=0$ for any
$w\in W\}$ for any subspace $W\subset V$. Recall
that $W$ is {\it isotropic} (or $\Phi$-\textit{isotropic}) if $W\subset W^\bot$.
\end{sub}

\begin{sub}{\bf Linear morphisms.}\label{subsection 2.2}
\rm

\begin{definition}\label{lin mor}
We call a morphism $\phi: X\to Y$ of algebraic varieties (or ind-varieties) \textit{linear} if $\phi$ induces an
epimorphism of Picard groups $\phi^*:\Pic Y\to\Pic X$.
\end{definition}

In this paper we focus on linear embeddings $\phi: X\to Y$ of grassmannians or isotropic grassmannians. In this case
$\phi$ is linear iff $\phi^*\mathcal{O}_Y(1)\cong\mathcal{O}_X(1)$.
By a \textit{projective space on}, or \textit{in}, a variety (or ind-variety) $X$ we
understand a linearly embedded subvariety $Y$ of $X$ isomorphic to a projective space.
Note that the Pl\"ucker embedding $G(k,V)\hookrightarrow\mathbb{P}(H^0(\mathcal{O}_{G(k,V)}(1))^*)$ is a linear
morphism.

By a \textit{quadric on} $X$ of dimension $m\ge3$ we understand a linearly embedded subvariety
$Y$ of $X$ isomorphic to a smooth $m$-dimensional quadric. By a quadric on $X$ of dimension 2
we understand the image of an embedding
$i:\mathbb{P}^1\times\mathbb{P}^1\hookrightarrow X$ such that $i^*\mathcal{O}_X(1)\simeq
\mathcal{O}_{\mathbb{P}^1}(1)\boxtimes\mathcal{O}_{\mathbb{P}^1}(1)$.
By a quadric on $X$ of dimension 1, or a \textit{conic} on $X$,
we understand the image of an embedding $i:\mathbb{P}^1\hookrightarrow X$ such that
$i^*\mathcal{O}_X(1)\simeq\mathcal{O}_{\mathbb{P}^1}(2)$. Given a quadric $Q$, we set $\mathbb{P}_Q=\mathbb{P}
(H^0(\mathcal{O}_Q(1))^*)$ for $m\ge3$, and respectively
$\mathbb{P}_Q=\mathbb{P}(H^0(\mathcal{O}_{\mathbb{P}^1}(1)\boxtimes\mathcal{O}_{\mathbb{P}^1}(1))^*)$,
$\mathbb{P}_Q=\mathbb{P}(H^0(\mathcal{O}_{Q}(2))^*)$ for $m=2,1$. Then $Q$ is canonically embedded into $\mathbb{P}_Q$.
\end{sub}


\begin{sub}{\bf Orthogonal grassmannians.}\label{linear ortho ind-grassm}
\rm

Let $\Phi\in S^2V^*$ be a non-degenerate symmetric form on $V$.
For $\dim V\ge3$ and  $1\le k\le[\frac{\dim V}{2}]$, the \textit{orthogonal grassmannian} $GO(k,V)$ is defined as the
subvariety of $G(k,V)$ consisting of $\Phi$-isotropic $k$-dimensional subspaces of $V$.
Unless $\dim V=2n$, $k=n$, $GO(k,V)$ is a smooth irreducible variety.
For $\dim V=2n$, $k=n$, $GO(k,V)$ is smooth and has two irreducible components, both of which are isomorphic to
$GO(n-1,V')$ where $\dim V'=2n-1$.

The orthogonal grassmannian $GO(k,V)$ has the following dimension:
$$
\dim GO(k,V)=\Bigl\{
\begin{array}{lll}
2kn-\frac{1}{2}k(3k+1) & \textrm{for}\ \ 1\le k\le n, & \dim V=2n, \\
k(2n+1)-\frac{1}{2}k(3k+1) & \textrm{for}\ \ 1\le k\le n, & \dim V=2n+1.
\end{array}
$$
Moreover, for any $V$ and $1\le k\le[\frac{\dim V}{2}]$, $k\ne\frac{\dim V}{2}-1$,
$$
\Pic GO(k,V)=\mathbb{Z}\mathcal{O}_{GO(k,V)}(1),
$$
where the sheaf $\mathcal{O}_{GO(k,V)}(1)$ posesses the following property: if $t: GO(k,V)\hookrightarrow G(k,V)$
is the tautological embedding, then
$$
t^*\mathcal{O}_{G(k,V)}(1)\cong\Bigl\{
\begin{array}{lll}
\mathcal{O}_{GO(k,V)}(1) & for & 1\le k\le [\frac{\dim V}{2}]-1, \\
\mathcal{O}_{GO(k,V)}(2) &  for & k=[\frac{\dim V}{2}].
\end{array}
$$
In what follows we will think of $GO(n-1,V)$ for $\dim V=2n$ as a variety of isotropic flags rather than as an
orthogonal grassmannian. In addition, we exclude the case $\dim V=2n,\ k=n$ from consideration.
More precisely, when writing $GO(k,V)$ below
we assume that $\dim V\ge7$ and $k\ne\frac{\dim V}{2}$, $k\ne\frac{\dim V}{2}-1$.

For $k<n=[\frac{\dim V}{2}]$ on $GO(k,V)$ there is a single family of maximal projective spaces of dimension $k$
with base $PO_{\alpha}(k,V)$. There is also a family of $(\dim V-2k)$-dimensional maximal quadrics not contained
in projective spaces on $GO(k,V)$. We denote the base of this family by $QO_{\beta}(k,V)$.
In addition, for $k\le[\frac{\dim V}{2}]-2$ there is a family of 4-dimensional maximal quadrics not contained
in projective spaces on $GO(k,V)$. We denote the base of this family by $QO_{\gamma}(k,V)$.

For $k=n$ on $GO(k,V)$ there is a single family of maximal projective spaces of
dimension $[\frac{\dim V-1}{2}]$ with irreducible base $PO_{\alpha}(k,V)$. Furthermore, if $\dim V=2n+1$ and
$1\le k\le n$, on $GO(k,V)$ there is a single family of maximal projective spaces of dimension $n-k$ with irreducible
base $PO_{\beta}(k,V)$.
The varieties $PO_{\alpha}(k,V)$, $PO_{\beta}(k,V)$, $QO_{\beta}(k,V)$ and $QO_{\gamma}(k,V)$ are described by
the following lemma.

\begin{lemma}\label{Bkn}
(i) If $1\le k\le n-1$, then each $k$-dimensional projective space on $GO(k,V)$ is of the form
\begin{equation}\label{PkI}
\{V_k\in GO(k,V)|\ V_k\subset V_{k+1}\}\simeq\mathbb{P}(V_{k+1}^*)
\end{equation}
for a fixed $(k+1)$-dimensional isotropic subspace $V_{k+1}$. Consequently, for $k\ne\frac{\dim V}{2}-2$,
$PO_{\alpha}(k,V)$ is isomorphic to $GO(k+1,V)$.

(ii) If $1\le k\le n-1$, then $PO_{\beta}(k,V)$ is isomorphic to the variety of isotropic $(k-1,n)$-flags in $V$,
and for any point $(V_{k-1}\subset V_n)\in PO_{\beta}(k,V)$ the corresponding projective space on $GO(k,V)$ is
\begin{equation}\label{Pn-k}
\{V_k\in GO(k,V)|\ V_{k-1}\subset V_k\subset V_n\}\simeq\mathbb{P}(V_n/V_{k-1}).
\end{equation}

(iii) If $k=n$, then $PO_{\alpha}(n,V)$ is isomorphic to $GO(n,V)$,
and for any point $V_n\in GO(n,V)$ the corresponding projective space on $GO(n,V)$ is
\begin{equation}\label{PnI}
\{V'_n\in GO(n,V)|\ \dim(V'_n\cap V_n)=n-1\}\simeq\mathbb{P}(V/V_n).
\end{equation}

(iv) If $1\le k\le n$, then $QO_{\beta}(k,V)$ is isomorphic to $GO(k-1,V)$, and for any point
$V_{k-1}\in GO(k-1,V)$ the corresponding quadric on $GO(k,V)$ is
\begin{equation}\label{Q}
\{V_k\in GO(k,V)|V_k\supset V_{k-1}\}\simeq GO(1,V_{k-1}^\bot/V_{k-1}).
\end{equation}

(v) $QO_{\gamma}(k,V)$ is isomorphic to the variety of isotropic $(k-2,k+2)$-flags in $V$, and for any point
$(V_{k-2}\subset V_{k+2})\in QO_{\gamma}(k,V)$ the corresponding quadric on $GO(k,V)$ is
\begin{equation}\label{Q1}
\{V_k\in GO(k,V)|V_{k-2}\subset V_k\subset V_{k+2}\}.
\end{equation}

(vi) Any maximal quadric on $GO(k,V)$ is either of the form (\ref{Q}) or (\ref{Q1}), or lies in a projective space on
$GO(k,V)$.
\end{lemma}

\begin{proof}
We leave the proof of (i)-(v) to the reader and give an outline of the proof of (vi).
Let $Q$ be a quadric on $GO(k,V)$ and let $G$ be the variety of projective planes in $\mathbb{P}_Q$.
In $G$ there is a dense open subset $U=\{\mathbb{P}^2\in G\ |\ \mathbb{P}^2\cap Q$ is a conic$\}$, and if
$\mathbb{P}^2\cap Q=C$ then $\mathbb{P}_C=\mathbb{P}^2\in U$.
In what follows, by a slight abuse of notation, we will indicate this latter fact by writing $C\in U$.


Let $F$ be the variety of $(1,k)$-isotropic flags in $V$
with projections $\mathbb{P}(V)\overset{pr_1}\leftarrow F\overset{pr_2}\to GO(k,V)$. For any $C\in U$ set
$\tilde{K}_C:=pr_2^{-1}(C)$, $K_C:=pr_1(\tilde{K}_C)$ and let $p_C:=pr_1|_{\tilde{K}_C}:\ \tilde{K}_C\to K_C$ be the
projection. There are three possibilities:

(a) there exists a dense open subset $U'$ in $U$ such that, for any $C\in U'$, $p_C$ is an isomorphism and $K_C$ is a
quadratic cone with vertex $S=\mathbb{P}(V_{k-1}(C))$ for some subspace $V_{k-1}(C)$ in $V$,

(b) there exists a dense open subset $U'$ in $U$ such that, for any $C\in U'$, $p_C$ is an isomorphism and $K_C$ is a
quadratic cone with vertex $S=\mathbb{P}(V_{k-2}(C))$ for some subspace $V_{k-2}(C)$ in $V$,

(c) for any $C\in U$, $p_C$ is a double covering and $K_C=\mathbb{P}(V_{k+1}(C))$ for some subspace $V_{k+1}(C)$ of $V$.

Using the fact that $U$ and $U'$ are dense subsets in $G$, one easily checks the following facts.
In case (a) the space $V_{k-1}=V_{k-1}(C)$ does not depend on the conic $C\in U'$ and
$Q\subset GO(1,V_{k-1}^\bot/V_{k-1})\in QO_{\beta}(k,V)$. In case (b)  the space $V_{k-2}=V_{k-2}(C)$ does not depend on
the conic $C\in U'$ and $Q$ is contained in a quadric $\bar{Q}$ given by formula (\ref{Q1}), i.e
$\bar{Q}\in QO_{\gamma}(k,V)$. In case (c) the space $V_{k+1}=V_{k+1}(C)$ does not depend on the conic $C\in U$, so that
$Q\subset\mathbb{P}(V_{k+1}^*)\subset GO(k,V)$.
\end{proof}

In what follows we will sometimes write $\mathbb{P}^k_{\alpha}$ for a maximal projective space on $GO(k,V)$ of the form
(\ref{PkI}) or (\ref{PnI}), and $\mathbb{P}^{n-k}_{\beta}$ for a maximal projective space on $GO(k,V)$ of the
form (\ref{Pn-k}). We will also write $Q^{\dim V-2k}_{\beta}$ for a maximal quadric on $GO(k,V)$ of the
form (\ref{Q}), and $Q^4_{\gamma}$ for a maximal quadric of the form (\ref{Q1}).


\begin{lemma}\label{two possib} Let $1\le k\le n$.

(i) The intersection of any two distinct projective spaces $\mathbb{P}^k_{\alpha}$ and
$(\mathbb{P}^k_{\alpha})'$ on $GO(k,V)$
is either empty or equals a point. The intersection of two distinct projective spaces $\mathbb{P}^{n-k}_{\beta}$ and $(\mathbb{P}^{n-k}_{\beta})'$ on $GO(k,V)$ can be any projective space of dimension less or equal $n-k-1$ or be empty.

(ii) The intersection of any projective space $\mathbb{P}^k_{\alpha}$ and any quadric
$Q^{\dim V-2k}_{\beta}$ on $GO(k,V)$ is empty, equals a point, or equals a projective line. The intersection of any two
distinct quadrics $Q^{\dim V-2k}_{\beta}$ and $(Q^{\dim V-2k}_{\beta})'$ on $GO(k,V)$ is either empty or equals a point.



(iii) Assume $k\le n-1$. Then the intersection of any two distinct projective spaces $\mathbb{P}^k_{\alpha}$ and $\mathbb{P}^{n-k}_{\beta}$
on $GO(k,V)$ is empty, equals a point, or equals a projective line.

(iv) Assume $k\le n-1$. Then $\mathbb{P}^k_{\alpha}\cap \mathbb{P}^{n-k}_{\beta}=\{V_k\}$ if and only if\
$\mathbb{P}^k_{\alpha}=\mathbb{P}(V_{k+1}^*)$, $\mathbb{P}^{n-k}_{\beta}=\mathbb{P}(V_n/V_{k-1})$ for a configuration
of isotropic vector subspaces $V_{k+1},V_{k-1},V_n$ of $V$ satisfying
$$
V_{k-1}\subset V_k\subset V_n,\ \ \ V_{k+1}\cap V_n=V_k.
$$
\end{lemma}

\begin{proof} Exercise.
\end{proof}


\begin{lemma}\label{family CO}
(i) Let $\mathbb{P}^1$ be a projective line on $GO(k,V)$, $x\not\in\mathbb{P}^1$ be a fixed
point in $GO(k,V)$, and $C\subset GO(k-1,V)$ be an irreducible curve such that, for any $V_{k-1}\in C$, the quadric
$GO(1,V_{k-1}^\bot/V_{k-1})$ on $GO(k,V)$ contains $x$ and intersects $\mathbb{P}^1$. Then $C$
is a projective line on $GO(k-1,V)$.

(ii) Assume $1\le k\le n-1$. Let $\mathbb{P}^1$ be a projective
line on $GO(k,V)$, $x\not\in\mathbb{P}^1$ be a fixed point in $GO(k,V)$, and $C\subset GO(k+1,V)$ be an
irreducible curve such that, for any $V_{k+1}\in C$, the projective space $\mathbb{P}(V_{k+1}^*)$ on $GO(k,V)$ contains
$x$ and intersects $\mathbb{P}^1$. Then $C$ is a projective line on $GO(k+1,V)$.
\end{lemma}

\begin{proof}
(i) Assume $k<n$ and let
\begin{equation}\label{pensil}
\mathbb{P}^1=\{V_k\in V|U_{k-1}\subset V_k\subset U_{k+1}\}
\end{equation}
for a fixed isotropic flag $U_{k-1}\subset U_{k+1}$ in $V$.
Next, let $x=W_k$. Since for any $V_{k-1}\in C$, the quadric $GO(1,V_{k-1}^\bot/V_{k-1})$ contains
the point $x$, we have $V_{k-1}\subset W_k$, and consequently
$$
\Span(\underset{V_{k-1}\in C}{\cup}V_{k-1})=W_k.
$$
The condition that the quadric $GO(1,V_{k-1}^\bot/V_{k-1})$ intersects $\mathbb{P}^1$ shows that
\begin{equation}\label{incl 1}
V_{k-1}\subset V_k\subset U_{k+1},\ \ \ U_{k-1}\subset V_k
\end{equation}
for some $V_k\in\mathbb{P}^1$. In particular,
$$
W_k\subset U_{k+1}.
$$
Note that $U_{k-1}\not\subset W_k$ as otherwise $x\in\mathbb{P}^1$.
Therefore $W_{k-2}:=W_k\cap U_{k-1}$ is a
$(k-2)$-dimensional subspace of $W_k$. Now (\ref{incl 1}) implies that
$C=\{V_{k-1}\in GO(k-1,V)|\ W_{k-2}\subset V_{k-1}\subset W_k\}$, i.e. $C$ is a projective line on $GO(k-1,V)$.

We leave the case $k=n$ to the reader.


(ii) Formula (\ref{pensil}) holds also in this case. Furthermore,
$$
\underset{V_{k+1}\in C}{\cap}V_{k+1}=W_k=x.
$$
For any $V_{k+1}\in C$, the condition that  $\mathbb{P}(V_{k+1}^*)$ intersects $\mathbb{P}^1$ yields
$V_k$ such that
$$
U_{k-1}\subset V_k\subset V_{k+1}.
$$
Therefore $U_{k-1}\subset W_k$. Now if $U_{k+1}\in C$, then for any $V_{k+1}\in C$, $\mathbb{P}(V_{k+1}^*)$ intersects
$\mathbb{P}^1$ in $x$, contrary to the assumption that $x\not\in\mathbb{P}^1$. Hence, $U_{k+1}\not\in C$ and one checks
that $C=\{V_{k+1}\subset V|W_k\subset V_{k+1}\subset W_{k+2}\}$, where
$W_{k+2}:=\Span(W_k,U_{k+1})$ is a $(k+2)$-dimensional subspace of $V$. This means
that $C$ is a projective line on $GO(k+1,V)$.


\end{proof}




\end{sub}




\begin{sub}{\bf Symplectic grassmannians.}\label{symp ind-grassm}

\rm
Let now $\Phi\in \wedge^2V^*$ be a non-degenerate symplectic form on $V$, $\dim V=2n$.

Assume $1\le k\le n$. Recall that the \textit{$k$-th symplectic grassmannian}  $GS(k,V)$ is the smooth irreducible
subvariety of $G(k,V)$ consisting of $\Phi$-isotropic $k$-dimensional subspaces of $V$.
It is well known that
\begin{equation}\label{dim IIkn}
\dim GS(k,V)=2kn-\frac{1}{2}k(3k-1).
\end{equation}
It is also known that $\Pic GS(k,V)=\mathbb{Z}\mathcal{O}_{GS(k,V)}(1)$ and
$\mathcal{O}_{GS(k,V)}(1)=i^*\mathcal{O}_{G(k,V)}(1)$, where $i:GS(k,V)\hookrightarrow G(k,V)$
is the tautological embedding.



\vspace{2mm}

One can see that, for $1\le k\le n-1$, there are two families of maximal
projective spaces on $GS(k,V)$ of respective dimensions $k$ and $2n-2k+1$, with bases
$PS_{\alpha}(k,V)$ and $PS_{\beta}(k,V)$.
For $k=n$ there is a single family $PS_{\beta}(n,V)$ of maximal projective lines on $GS(k,V)$.
\begin{lemma}\label{Bkn sympl}

(i) Let $1\le k\le n-1$. Then $PS_{\alpha}(k,V)$ is isomorphic to $GS(k+1,V)$, and for any point $V_{k+1}\in GS(k+1,V)$
the corresponding projective space on $GS(k,V)$ is
\begin{equation}\label{PkS}
\{V_k\in GS(k,V)|\ V_k\subset V_{k+1}\}\simeq\mathbb{P}(V_{k+1}^*).
\end{equation}

(ii) Let $1\le k\le n$. Then $PS_{\beta}(k,V)$ is isomorphic to $GS(k-1,V)$, and for any point
$V_{k-1}\in GS(k-1,V)$ the corresponding projective space on $GS(k,V)$ is
\begin{equation}\label{Z(V_k)}
\{V_k\in GS(k,V)\ |\ V_{k-1}\subset V_k\subset V_{k-1}^\bot\}\simeq\mathbb{P}(V_{k-1}^\bot/V_{k-1}).
\end{equation}

(iii) If $k=n$, then any maximal projective space on $GS(n,V)$ is a projective line.
\end{lemma}
\begin{proof} Exercise.
\end{proof}

In what follows we will sometimes write $\mathbb{P}^k_{\alpha}$ for a maximal projective space on $GS(k,V)$ of the form
(\ref{PkS}), and $\mathbb{P}^{2n-2k+1}_{\beta}$ for a maximal projective space on $GS(k,V)$ of the
form (\ref{Z(V_k)}) (despite the fact that we use the same notation as in the orthogonal case,
we will carefully distinguish between the two cases).



\begin{lemma}\label{intersect dim=1}
Let $\dim V=2n$, $n\ge2$, and $1\le k\le n-1$.

(i) The intersection of any two distinct projective spaces $\mathbb{P}^k_{\alpha}$ and
$(\mathbb{P}^k_{\alpha})'$ (respectively, $\mathbb{P}^{2n-2k+1}_{\beta}$ and $(\mathbb{P}^{2n-2k+1}_{\beta})'$) on
$GS(k,V)$ is either empty or equals a point.

(ii) The intersection of any two distinct projective spaces $\mathbb{P}^k_{\alpha}$ and
$\mathbb{P}^{2n-2k+1}_{\beta}$ on $GS(k,V)$ is either empty or equals a projective line.

\vspace{3mm}
(iii) The spaces $\mathbb{P}^k_{\alpha}$ and $\mathbb{P}^{2n-2k+1}_{\beta}$ intersect in a projective line
if and only if
$\mathbb{P}^k_{\alpha}=\mathbb{P}(V_{k+1}^*)$, $\mathbb{P}^{2n-2k+1}_{\beta}=\mathbb{P}(V_{k-1}^\bot/V_{k-1})$ for a
flag $V_{k-1}\subset V_{k+1}$ of isotropic subspaces of $V$. Then
$\mathbb{P}^k_{\alpha}\cap\mathbb{P}^{2n-2k+1}_{\beta}=\mathbb{P}(V_{k+1}/V_{k-1})$.
\end{lemma}
\begin{proof} Exercise.
\end{proof}

\begin{lemma}\label{family CS}
(i) Assume $2\le k\le n$.
Let $\mathbb{P}^1$ be a projective
line on $GS(k,V)$, $x\not\in\mathbb{P}^1$ be a fixed point in $GS(k,V)$, and $C\subset GS(k-1,V)$ be
an irreducible curve such that, for any $V_{k-1}\in C$, the projective space $\mathbb{P}(V_{k-1}^\bot/V_{k-1})$ on
$GS(k,V)$ contains $x$ and intersects $\mathbb{P}^1$. Then $C$ is a projective line on $GS(k-1,V)$.

(ii) Assume $1\le k\le n-1$. Let $\mathbb{P}^1$ be a projective
line on $GS(k,V)$, $x\not\in\mathbb{P}^1$ be a fixed point in $GS(k,V)$, and $C\subset GS(k+1,V)$ be an
irreducible curve such that, for any $V_{k+1}\in C$, the projective space $\mathbb{P}(V_{k+1}^*)$ on $GS(k,V)$ contains
$x$ and intersects $\mathbb{P}^1$. Then $C$ is a projective line on $GS(k+1,V)$.


\end{lemma}
\begin{proof}
Very similar to the proof of Lemma \ref{family CO}.
\end{proof}

\vspace{3mm}




\end{sub}





\vspace{1cm}
\section{Linear embeddings of grassmannians}\label{Lin emb}

\vspace{1cm}

In this section we study linear embeddings of grassmannians and isotropic grassmannians.

We start with the following general lemma whose proof we leave to the reader.


\begin{lemma}\label{finite}
Any non-constant morphism of grassmannians (respectively, orthogonal or symplectic grassmannians) is finite.
\end{lemma}


\begin{definition}\label{standard gen}
Let $X,X'$ be grassmannians. An embedding $\phi':X\hookrightarrow X'$
is a \textit{standard extension}, if there are isomorphisms $i_X,i_{X'}$ and an
embedding $\phi:G(k,V)\hookrightarrow G(k',V')$ for $\dim V'\ge\dim V,\ k'\ge k,$ such
that the diagram
\begin{equation}\label{Def 3.1}
\xymatrix{
X \ar@{^{(}->}[r]^-{\phi'}\ar[d]^-{i_X} & X'\ar[d]^-{i_{X'}}\\
G(k,V) \ar@{^{(}->}[r]^-{\phi} & G(k',V') }
\end{equation}
is commutative and $\phi$ is given by the formula
\begin{equation}\label{st ext}
 \phi:V_k\mapsto V_k\oplus W
\end{equation}
for some fixed isomorphism $V'\simeq V\oplus\hat{W}$ and a fixed subspace $W\subset\hat{W}$ of dimension $k'-k$.
\end{definition}


It is easy to see that a standard extension is a linear embedding. Furthermore, if $\mathbb{P}^q$ is a projective space
on $G(k,V)$, then the inclusion $\mathbb{P}^q\hookrightarrow G(k,V)$ is a standard extension.

\begin{example}
Let $V'=V\oplus\hat{W},\ X=G(n-k,V^*),\ X'=G(k',V'),\ W\subset\hat{W}$ be a fixed subspace of dimension $k'-k$,
$\varepsilon: V\to V$ be any autorphism. Then the morphism
$$
X=G(n-k,V^*)\simeq G(k,V)\to G(k',V')=X',
$$
$$
V_{n-k}\mapsto V_{n-k}^*\mapsto\varepsilon(V_{n-k}^*)\oplus W
$$
is a standard extension. Here $i_X$ is the isomorphism $G(n-k,V^*)\simeq G(k,V)$ and $i_{X'}$ is the automorphism
of $X'$ induced by the automorphism $\varepsilon^{-1}\oplus{\rm id}_{\hat{W}}$ of $V'$.
\end{example}


\begin{remark}\label{intrinsic def}
Note that, for a standard extension $\phi':X\to X'$ the dimensions of $V$ and $V'$ are fixed by the respective
isomorphism classes of $X$ and $X'$, however, the choice between $k$ and $\dim V-k$, respectively, $k'$ and
$\dim V'-k'$, in diagram (\ref{Def 3.1}) is made by the morphism $\phi'$.
Furthermore, if $V,V',k,k'$ are chosen, fixing a standard extension $\phi:G(k,V)\to G(k',V')$
for which $i_{G(k,V)}$ and $i_{G(k',V')}$ are automorphisms is equivalent to fixing some linear algebraic data.
More precisely, given such a standard extension $\phi:G(k,V)\hookrightarrow G(k',V'),$ we can recover $W$ by the formula
$W=\underset{V_k\in G(k,V)}\cap \phi(V_k)$. Set $U:=\Span(\underset{V_k\in G(k,V)}\cup \phi(V_k))$.
Then $W\subset U$ is a flag in $V'$ and $\phi$ determines a surjective linear operator
$\underline{\phi}:U\to V$ with kernel $W$, such that $(\underline{i})^{-1}(V_k)=\phi(V_k)$
for any $k$-dimensional subspace $V_k\in G(k,V)$. It is easy to check that fixing the standard extension $\phi$ is
equivalent to fixing the triple $(W,U,\underline{\phi})$.

In what follows we will write somewhat informally $\phi:G(k,V)\hookrightarrow G(k',V')$ for a general standard
extension, while we will speak about a \textit{strict} standard extension when $i_{G(k,V)}$ and $i_{G(k',V')}$ are
automorphisms. Given a strict standard extension $\phi:G(k,V)\hookrightarrow G(k',V')$, the isomorphism
$V'\simeq V\oplus\hat{W}$ can always be changed so that $\phi$ is given simply by formula (\ref{st ext}).

\end{remark}

We now give a similar definition of a standard extension of isotropic grassmannians (cf. [DP] and [PT1, section 3]).
\begin{definition}\label{lin mor isotr}
An embedding $\phi:GO(k,V)\hookrightarrow GO(k',V')$ is a \textit{standard extension} if
$\phi$ is given by formula (\ref{st ext}) for some orthogonal isomorphism
$V'\simeq V\oplus\hat{W}$ and a fixed isotropic subspace $W$ of $\hat{W}$.
A standard extension of symplectic grassmannians is defined in the same way by replacing $GO$ with $GS$,
and the orthogonal isomorphism $V'\simeq V\oplus\hat{W}$ by a symplectic isomorphism $V'\simeq V\oplus\hat{W}$.
\end{definition}

Under an orthogonal isomorphism (respectively, symplectic isomorphism) we mean an isomorphism of vector spaces together
with an isomorphism of forms $\Phi'\simeq\Phi\oplus\hat{\Phi}$, where $\Phi$ is a fixed symmetric (respectively,
symplectic) form on $V$, $\Phi'$ is a fixed (respectively, symplectic) form on $V'$, and $\hat{\Phi}$ is a fixed
symmetric (respectively, symplectic) form on $\hat{W}$.

\begin{remark}\label{intrinsic def2}
A standard extension of isotropic grassmannians can be defined as follows: consider a flag of subspaces
$W\subset U$ of $V'$, where $W$ is isotropic and there is a surjective linear operator $\underline{\phi}:U\to V$ with
kernel $W$, such that the form $\underline{\phi}^*\Phi$ coincides with the form induced on $U$ by the form $\Phi'$.
This datum defines an embedding $GO(k,V)\to GO(k',V')$ (respectively, $GS(k,V)\to GS(k',V')$) by the formula
$$
\phi:\ V_k\mapsto(\underline{\phi})^{-1}(V_k)\subset U\subset V'\ \ \ \textrm{for} \ \ \ V_k\in GO(k,V) \ \ \
(\textrm{resp.,}\ V_k\in GS(k,V)).
$$
Furthermore,
\begin{equation}\label{gen W,U}
W=\cap \phi(V_k),\ \ \ U=\Span(\cup \phi(V_k)),
\end{equation}
where $V_k$ runs over $GO(k,V)$ (respectively, $GS(k,V)$) and the intersection and the union are taken in $V'$.
\end{remark}


\begin{remark}\label{prop of st ext1}
Let $\phi:G(k,V)\to G(k',V')$ be a strict standard extension (respectively, $\phi:GO(k,V)\to GO(k',V')$ or
$\phi:GS(k,V)\to GS(k',V')$ be a standard extension). Then
\begin{equation}\label{stand ineq1}
k'\ge k\ \ \  \textrm{and}\ \ \ \ \dim V'-k'\ge\dim V-k\ge0
\end{equation}
(respectively,
\begin{equation}\label{stand ineq2}
k'\ge k\ \ \  \textrm{and}\ \ \ \frac{1}{2}\dim V'-k'\ge\frac{1}{2}\dim V-k\ge0).
\end{equation}
Indeed, Definition \ref{standard gen} implies $k'-k=\dim W\ge0$.
Next, from $\dim W\le\dim\hat{W}=\dim V'-\dim V$
it follows that $\dim V'-k'=\dim V-k+(\dim\hat{W}-\dim W)\ge\dim V-k$. This proves (\ref{stand ineq1}).
As for (\ref{stand ineq2}), from Definition
\ref{lin mor isotr} we have $k'-k=\dim W\ge0$. Furthermore, as $V_k$ is $\Phi$-isotropic,
$V_{k'}:=V_k\oplus W$ is $\Phi'$-isotropic and $W$ is $\hat{\Phi}$-isotropic, we have
$k\le\frac{1}{2}\dim V,\ k'\le\frac{1}{2}\dim V',\ 0\le\dim W\le\frac{1}{2}\dim\hat{W}=\frac{1}{2}(\dim V'-\dim V)$.
This implies $\frac{1}{2}\dim V'-k'=\frac{1}{2}\dim V-k+\frac{1}{2}\dim\hat{W}-\dim W\ge\frac{1}{2}\dim V-k\ge0$.
\end{remark}

\begin{definition}\label{isotr st ext}
(a) Let $V''$ be an isotropic subspace of $V$. For $\mathbb{Z}_+\ni k\le\dim V''$, we call
the natural inclusions $G(k,V'')\hookrightarrow GO(k,V)$ and $G(\dim V''-k,{V''}^*)\hookrightarrow GO(k,V)$
(respectively, $G(k,V'')\hookrightarrow GS(k,V)$ and $G(\dim V''-k,{V''}^*)\hookrightarrow GS(k,V)$)
\textit{isotropic extensions}.

(b) A \textit{combination of isotropic and standard extensions} is an embedding of the form
$$
GO(k,V)\overset{t}\hookrightarrow G(k,V)\overset{\phi'}\hookrightarrow G(l,U)\overset{\tau}\hookrightarrow
GO(l,\tilde{U})\overset{\phi''}\hookrightarrow GO(k',V')
$$
(respectively,
$$
GS(k,V)\overset{t}\hookrightarrow G(k,V)
\overset{\phi'}\hookrightarrow G(l,U)\overset{\tau}\hookrightarrow GS(l,\tilde{U})\overset{\phi''}\hookrightarrow
GS(k',V')),
$$
where $t$ is the tautological embedding, $\phi'$ and $\phi''$ are standard extensions and $\tau$ is an
isotropic extension.
\end{definition}

Note that a combination of isotropic and standard extensions is always given by one of the formulas
$V_k\mapsto V_k\oplus W$ or
$V_k\mapsto V_k^\bot\oplus W$ for an appropriately chosen orthogonal (respectively, symplectic) isomorphism
$V'\simeq V\oplus\hat{W}$ and an isotropic subspace $W\subset\hat{W}$. Here $\bot$ refers to the orthogonal
(respectively, symplectic) structure on $V$. Furthermore, one easily proves the following lemma.

\begin{lemma}\label{compos of combin}
A composition of combinations of isotropic and standard extensions is a combination of isotropic and standard
extensions.
\end{lemma}

\begin{remark}\label{prop of st ext}
Let $\phi':X\to X'$ be a standard extension, where $X$ and $X'$ are both grassmannians or, respectively, isotropic
grassmannians of the same type. It is easy to see that, if $X$ and $X'$ are not (isomorphic to) projective spaces,
then $\phi'$ does not factor through an embedding of a projective space into $X'$. If $X$ and $X'$ are isotropic, then
$\phi'$ is not a combination of isotropic and standard extensions.
\end{remark}








{\bf Theorem 1}.
\textit{Let $X\simeq G(k,V)$, $X'\simeq G(k',V')$, or $X=GO(k,V),\ X'=GO(k',V')$, or
$X=GS(k,V),\ X'=GS(k',V')$, and let $\phi: X \to X'$ be a linear morphism. If $X=GO(k,V),\ X'=GO(k',V')$, assume
in addition that either $k\le[\frac{\dim V}{2}]-3$ and $k'\le[\frac{\dim V'}{2}]-3$, or that
$[\frac{\dim V'}{2}]-k'\le[\frac{\dim V}{2}]-k\le2$ and both $\dim V$ and $\dim V'$ are odd.
Then some of the following statements holds:\\
(i) $\phi$ is a standard extension;\\
(ii) $X$ and $X'$ are isotropic grassmannians and $\phi$ is a combination of isotropic and standard
extensions;\\
(iii) $\phi$ factors through a projective space on $X'$ or, in case $X'=GO(k',V')$, through a maximal quadric
$Q_{\beta}^{\dim V'-2k'}$.}

\begin{proof}
We first consider in detail the case of symplectic grassmannians.
The proof goes by induction on $k$. For $k=1$ the symplectic grassmannian
$GS(1,V)$ equals $\mathbb{P}(V)$, hence the linear morphism $\phi$ maps it isomorphically onto a projective space in
$X'$. Therefore statement (iii) holds trivially in this case.

Assume now that $k\ge2$ and the assertion holds for $k-1$ and any $k'\ge1$.
Set $n:=\frac{1}{2}\dim V$, $n':=\frac{1}{2}\dim V'$, $Y_{\beta}:=GS(k-1,V)$.
Let $Z:=\{(V_{k-1},x)\in Y_{\beta}\times X|\ x\in\mathbb{P}(V_{k-1}^\bot/V_{k-1})\}\overset{p}\to Y_{\beta}$
be the family of projective spaces $\mathbb{P}^q_\beta,\ q=2n-2k+2,$ on $X\simeq GS(k,V)$.
Since $\phi$ is a linear morphism, for any $V_{k-1}\in Y_{\beta}$ the morphism
$\phi(\mathbb{P}(V_{k-1}^\bot/V_{k-1}))$ is a projective space on $X'$.
Therefore we obtain a family
$\tilde{Z}:=\{(V_{k-1},x)\in Y_{\beta}\times X'|\ x\in\phi(\mathbb{P}(V_{k-1}^\bot/V_{k-1}))\}
\overset{\tilde{p}}\to Y_{\beta}$ of $q$-dimensional projective spaces on $X'$.
We claim that

(a) all spaces of the family $\tilde{p}:\tilde{Z}\to Y_{\beta}$ lie in the spaces of the
family with base $Y'_{\alpha}:=GS(k'+1,V')$ (this is possible only if $k\le n-1$),\\
or

(b) all spaces of the family
$\tilde{p}:\tilde{Z}\to Y_{\beta}$ lie in the spaces of the family with base $Y'_{\beta}:=GS(k'-1,V')$.


Indeed, consider the varieties $\Sigma_{\alpha}:=\{(V_{k-1},V'_{k'+1})\in Y_{\beta}\times Y'_{\alpha}|\
\phi(\mathbb{P}(V_{k-1}^\bot/V_{k-1}))\subset\mathbb{P}((V'_{k'+1})^*)\}$
and $\Sigma_{\beta}:=\{(V_{k-1},V'_{k'-1})\in Y_{\beta}\times Y'_{\beta}|\ \phi(\mathbb{P}(V_{k-1}^\bot/V_{k-1})
\subset \mathbb{P}(V_{k'-1}^\bot/V_{k'-1})\}$
with natural projections
$Y_{\beta}\overset{p_\alpha}\leftarrow\Sigma_{\alpha}\overset{q_{\alpha}}\to Y'_{\alpha}$
and
$Y_{\beta}\overset{p_\beta}\leftarrow\Sigma_{\beta}\overset{q_{\beta}}\to Y'_{\beta}$.
By construction,
$\Sigma_{\alpha}$ is a closed subset of $Y_{\beta}\times Y'_{\alpha}$ and $p_{\alpha}$ is a projective morphism. Hence,
$W_{\alpha}:=p_{\alpha}(\Sigma_{\alpha})$ is a closed subset of $Y_{\beta}$. By a similar reason,
$W_{\beta}:=p_{\beta}(\Sigma_{\beta})$ is a closed subset of $Y_{\beta}$. Since any space of the family
$\tilde{Z}\to Y_{\beta}$ lies
in at least one maximal space on $X'$, it follows that $W_{\alpha}\cup W_{\beta}=Y_{\beta}$. However, $Y_{\beta}$ is
irreducible, therefore either $W_{\alpha}=Y_{\beta}$ (i.e. case (a) holds), or
$W_{\beta}=Y_{\beta}$ (i.e. case (b) holds).


We now consider the cases (a) and (b) separately.

In the case (a), by Lemma \ref{intersect dim=1},(i), each space of the family $\tilde{p}:\tilde{Z}\to Y_{\beta}$ lies
in a unique space $\mathbb{P}((V'_{k'+1})^*)$ of the family with base $Y'_{\alpha}$. This
means that $p_{\alpha}:\Sigma_{\alpha}\to Y_{\beta}$ is a bijective morphism, hence an isomorphism as $Y_{\beta}$ is a
smooth variety. Therefore, there is a well defined morphism
\begin{equation}\label{mor i}
\phi_{\alpha}:=q_{\alpha}\circ p_{\alpha}^{-1}:\ Y_{\beta}=GS(k-1,V)\to Y'_{\alpha}=GS(k'+1,V').
\end{equation}
Moreover, there is a commutative diagram
\begin{equation}\label{phi Y^2}
\xymatrix{
& \Gamma\ar[rrr]^{\phi_{\Gamma}}\ar[dl]_{p_1}\ar[dr]^{p_2} &&& \Gamma'_{\alpha}\ar[dl]_{\bar{p}_1}\ar[dr]^{\bar{p}_2}& \\
Y_{\beta}\ar[dr]&& X\ar[dr]& Y'_{\alpha} && X',\\
&\ar[r]_{\phi_{\alpha}}&\ar[ur]&\ar[r]_{\phi}&\ar[ur]& }
\end{equation}
where $\Gamma$ is the variety of isotropic $(k-1,k)$-flags in $V$,
$\Gamma'_{\alpha}$ is the variety of isotropic $(k',k'+1)$-flags in $V'$,
and $\phi_\Gamma,\ p_1,\ p_2,\ \bar{p}_1$ and $\bar{p}_2$ are the induced projections.

Assume that $\phi_{\alpha}$ is not a constant map. We first show that $\phi_{\alpha}$ is linear.
Fix $V_{k+1}\in GS(k+1,V)$ and a subspace $V_{k-2}$ of $V_{k+1}$.
Consider the
projective plane $\mathbb{P}^2_X:=\mathbb{P}((V_{k+1}/V_{k-2})^*)$ on $X$.
The points on $\mathbb{P}^2_X$ are $k$-dimensional subspaces
$U_k\subset V$ such that $V_{k-2}\subset U_k\subset V_{k+1}$.
According to Lemma \ref{Bkn sympl},(i), any $U_k$ defines a projective space $\mathbb{P}(U_k^*)$ on $Y_{\beta}$,
and also a projective line
$\mathbb{P}((U_k/V_{k-2})^*)$ on $Y_{\beta}$.
Fix $U_k$ and denote the projective line $\mathbb{P}((U_k/V_{k-2})^*)$ by $\mathbb{P}^1_{Y_{\beta}}$.
Furthermore, fix a projective line $\mathbb{P}^1_X$ in $\mathbb{P}^2_X$
and consider the rational curve $\mathbb{P}^1_{\Gamma}:=
\{(V_{k-1},V_k)\in\Gamma|\ V_{k-1}\in\mathbb{P}^1_{Y_{\beta}},\ V_k\in \mathbb{P}^1_X\}$ on $\Gamma$.
Diagram (\ref{phi Y^2}) yields a commutative diagram
\begin{equation}\label{curves}
\xymatrix{
\mathbb{P}^1_{Y_{\beta}}\ar[d]_-{\phi_{\alpha}|_{\mathbb{P}^1_{Y_{\beta}}}} &&
\mathbb{P}^1_{\Gamma}\ar[rr]^-{p_2|_{\mathbb{P}^1_{\Gamma}}}
\ar[d]_{\phi_{\Gamma}|_{\mathbb{P}^1_{\Gamma}}}\ar[ll]_-{p_1|_{\mathbb{P}^1_{\Gamma}}} &&
\mathbb{P}^1_X\ar[d]_{\phi|_{\mathbb{P}^1_X}}  \\
\phi_{\alpha}(\mathbb{P}^1_{Y_{\beta}}) &&
\phi_{\Gamma}(\mathbb{P}^1_{\Gamma})\ar[rr]^-{\bar{p}_2|_{\phi_{\Gamma}(\mathbb{P}^1_{\Gamma})}}
\ar[ll]_-{\bar{p}_1|_{\phi_{\Gamma}(\mathbb{P}^1_{\Gamma})}} && \phi(\mathbb{P}^1_X).}
\end{equation}
Since $\phi:X\to X'$ is linear, $\phi|_{\mathbb{P}^1_X}:\mathbb{P}^1_X\to \phi(\mathbb{P}^1_X)$ is an isomorphism.
Furthermore, $p_2|_{\mathbb{P}^1_{\Gamma}}$ is an isomorphism by construction. Therefore,
$\bar{p}_2|_{\phi_{\Gamma}(\mathbb{P}^1_{\Gamma})}$ and $\phi_{\Gamma}|_{\mathbb{P}^1_{\Gamma}}$ are isomorphisms.

We claim now that $p_1|_{\mathbb{P}^1_{\Gamma}}$ and $\bar{p}_1|_{\phi_{\Gamma}(\mathbb{P}^1_{\Gamma})}$ are also
isomorphisms. Indeed, for $p_1|_{\mathbb{P}^1_{\Gamma}}$ this holds by construction. Consider
$\bar{p}_1|_{\phi_{\Gamma}(\mathbb{P}^1_{\Gamma})}$. As $\phi(\mathbb{P}^1_X)$ is a projective line in $X'$, the
subspaces of $V'$ corresponding to the points of $\phi(\mathbb{P}^1_X)$ lie in some $k'+1$-dimensional subspace
$V'_{k'+1}$ of $V'$. This implies in view of Lemma \ref{intersect dim=1},(i) that, for any two distinct points
$V'_{k'},V''_{k'}\in\phi(\mathbb{P}^1_X)$, the projective spaces $\mathbb{P}({V'}_{k'}^\bot/V'_{k'})$ and
$\mathbb{P}({V''_{k'}}^\bot/V''_{k'})$ on $Y'_{\alpha}$ have $V'_{k'+1}$ as unique common point.
Note that, for each $V'_{k'}\in\phi(\mathbb{P}^1_X)$, $\mathbb{P}({V'_{k'}}^\bot/V'_{k'})$
is the isomorphic image under $\bar{p}_1$ of the projective space $\bar{p}_2^{-1}(V'_{k'})$, and that
$\bar{p}_2^{-1}(V'_{k'})\cap\phi_{\Gamma}(\mathbb{P}^1_{\Gamma})$ is a single point.
Hence, either $\bar{p}_1(\phi_{\Gamma}(\mathbb{P}^1_{\Gamma}))=\phi_{\alpha}(\mathbb{P}^1_{Y_{\beta}})$ equals the point
$V'_{k'+1}$, or $\bar{p}_1|_{\phi_{\Gamma}(\mathbb{P}^1_{\Gamma})}$ is an isomorphism. However, the former case is
impossible since $\phi_{\alpha}|_{\mathbb{P}^1_{Y_{\beta}}}$ is a non-constant, hence finite morphism by Lemma
\ref{finite}. Thus $\bar{p}_1|_{\phi_{\Gamma}(\mathbb{P}^1_{\Gamma})}$ is an isomorphism.

Diagram (\ref{curves}) implies now that $\phi_{\alpha}|_{\mathbb{P}^1_{Y_{\beta}}}$ is also an isomorphism.
To show that $\phi_{\alpha}$ is linear it suffices to prove that $\phi_{\alpha}(\mathbb{P}^1_{Y_{\beta}})$
is a projective line on $Y'_{\alpha}$. This latter fact follows directly from Lemma \ref{family CS},(ii)
applied to the following data: $\mathbb{P}^1=\mathbb{P}^1_{Y_{\beta}}$,
$x=\phi(\Span(\underset{V_{k-1}\in\mathbb{P}^1}\cup V_{k-1}))$, $C=\phi_{\alpha}(\mathbb{P}^1_{Y_{\beta}})$.



Note next that the diagram (\ref{phi Y^2}) allows to reconstruct $\phi$ from $\phi_{\alpha}$. Indeed, for any
$V_k\in X$ the projective space $\mathbb{P}(V_k^*)=p_1(p_2^{-1}(V_k))$ is mapped via
$\phi_{\alpha}$ to the unique projective space $\mathbb{P}({V'_{k'}}^\bot/V'_{k'})$ which contains
$\phi_{\alpha}(\mathbb{P}^{k-1}_{\alpha})$. The original morphism $\phi$ is precisely the map assigning $V'_{k'}$ to
$V_k\in X$.

We are now ready to apply the induction assumption to $\phi_{\alpha}$. Since $\phi_{\alpha}$ is linear we conclude that
there are the following three possibilities: (a.1) $\phi_{\alpha}$ is a standard extension, or (a.2) $\phi_{\alpha}$
factors through an isotropic extension, or (a.3) $\phi_{\alpha}$ factors through a morphism to a projective space in
$Y_\alpha$.

(a.1) In this case we have a fixed isomorphism $V'\simeq V\oplus \hat{W}$ and $\phi_{\alpha}$ is given by the formula
\begin{equation}\label{st eq 2}
V_{k-1}\mapsto V_{k-1}\oplus W
\end{equation}
for an isotropic subspace $W$ of $\hat{W}$ (see Remark \ref{intrinsic def}). Therefore, for any $V_k\in X$, the
space $\mathbb{P}^{k-1}_{\alpha}=\mathbb{P}(V_k^*)$ on $Y_{\beta}$ is embedded by $\phi_{\alpha}$ in the projective
space $\mathbb{P}^{k'+1}_{\alpha}=\mathbb{P}((V_k\oplus W)^*)$ on $Y'_{\alpha}$.
Since $\mathbb{P}(V_k^*)=p_1(p_2^{-1}(V_k))$, diagram (\ref{phi Y^2}) implies that
$\phi_{\alpha}(\mathbb{P}(V_k^*))\subset\mathbb{P}^{\dim V'-2k'-1}_{\beta}=\bar{p}_1(\bar{p}_2^{-1}(V'_{k'}))$, where
$V'_{k'}:=\phi(V_k)$. We have thus shown that $\phi_{\alpha}(\mathbb{P}(V_k^*))$ lies in the intersection of maximal
projective spaces from the distinct families $PS_{\alpha}(k'+1,V')$ and $PS_{\beta}(k'+1,V')$ on $Y'_{\alpha}$. Hence,
by Lemma \ref{intersect dim=1},(ii), $k-1=\dim\phi_{\alpha}(\mathbb{P}(V_k^*))\le1$, i.e. $k=2$. Therefore,
$X\simeq GS(2,V),\ Y_{\beta}=\mathbb{P}(V)$, and $\phi_{\alpha}$ is an embedding of $\mathbb{P}(V)$ into
$Y'_{\alpha}$. Then $\phi_{\alpha}(\mathbb{P}(V))$ lies in a unique maximal projective space
$\mathbb{P}((V'_{k'+2})^*)$ for an isotropic subspace $V'_{k'+2}$ of $V'$.
This yields a monomorphism $j:V\hookrightarrow(V'_{k'+2})^*$.
Now the above reconstruction of $\phi$ via $\phi_{\alpha}$ shows that $\phi$ decomposes as
\begin{equation}\label{decomp1}
X=GS(2,V)\overset{t}\hookrightarrow G(2,V)\overset{\tilde{j}}\hookrightarrow
G(2,(V'_{k'+2})^*)\simeq G(k',V'_{k'+2})\overset{\tau}\hookrightarrow GS(k',V')=X',
\end{equation}
where $t$ is the tautological
embedding, the embedding $\tilde{j}$ is induced by the monomorphism $j$, and the embedding $\tau$ is induced by the
embedding of $V'_{k'+2}$ in $V'$.
Hence, $\phi$ is a combination of isotropic and standard extensions. One checks that as a consequence $\phi_{\alpha}$
is also a combination of isotropic and standard extensions. However, this contradicts to Remark \ref{prop of st ext},
and we conclude that case (a.1) is impossible.

(a.2) In this case  $\phi_{\alpha}$ is given by one of the formulas
\begin{equation}\label{st eq 2a}
V_{k-1}\mapsto V_{k-1}\oplus W
\end{equation}
or
\begin{equation}\label{st eq 2b}
V_{k-1}\mapsto V^\bot_{k-1}\oplus W,
\end{equation}
where $\bot$ refers to the symplectic structure on $V$. If $\phi_{\alpha}$ is given by (\ref{st eq 2a}), then
for an arbitrary $V_k\in X$,
$\phi_{\alpha}(\mathbb{P}(V^*_k))\subset\mathbb{P}((V_k\oplus W)^*)$.
Assume that $\dim\phi_{\alpha}(\mathbb{P}(V^*_k))>1$. Then by Lemma \ref{intersect dim=1},(ii),
$\phi_{\alpha}(\mathbb{P}(V^*_k))\not\subset\mathbb{P}(\phi(V_k)^{\bot'}/\phi(V_k))$,
where the symbol $\bot'$ refers to the symplectic structure on $V'$.
On the other hand, in view of diagram
(\ref{phi Y^2}), $\phi_{\alpha}(\mathbb{P}(V^*_k))\subset\mathbb{P}(\phi(V_k)^{\bot'}/\phi(V_k))$.
This implies $k=2$. Therefore, $\phi_{\alpha}$ is a combination of isotropic and standard extensions of the form
$$
Y_{\beta}=\mathbb{P}(V)\overset{\phi'_{\alpha}}\hookrightarrow G(l,U)
\overset{\tau_{\alpha}}\hookrightarrow
GS(l,\tilde{U})\overset{\phi''_{\alpha}}\hookrightarrow GS(k'+1,V')=Y'_{\alpha}
$$
(see Definition \ref{isotr st ext}). Then using diagram (\ref{phi Y^2}) it is easy to check that $\phi$ is given by
the formula
$$
V_k\mapsto V_k\oplus W
$$
and is a combination of isotropic and standard extensions of the form
$$
X=GS(2,V)\overset{t}\hookrightarrow G(2,V)\overset{\phi'}\hookrightarrow G(l+1,U)
\overset{\tau}\hookrightarrow
GS(l+1,\tilde{U})\overset{\phi''}\hookrightarrow GS(k',V')=X'.
$$



If $\phi_{\alpha}$ is given by (\ref{st eq 2b}), then
for an arbitrary $V_k\in X$,
$\phi_{\alpha}(\mathbb{P}(V^*_k))=\mathbb{P}((V^\bot_k\oplus W)^{\bot'}/(V^\bot_k\oplus W))$.
In view of the diagram (\ref{phi Y^2}) $\phi$ is given in this case by the formula
$$
V_k\mapsto V^\bot_k\oplus W
$$
and is a combination of isotropic and standard extensions of the form
$$
X=GS(k,V)\overset{t}\hookrightarrow G(k,V)\overset{\bot}{\simeq}G(2n-k,V)\overset{\phi'}\hookrightarrow
G(l+1,U)\overset{\tau}\hookrightarrow GS(l+1,\tilde{U})\overset{\phi''}\hookrightarrow GS(k',V')=X'.
$$
In this way, (ii) holds under the assumption (a.2).


(a.3) Here $\phi_{\alpha}$ factors through a morphism to some projective space $\mathbb{P}^s$ in $Y'_{\alpha}$,
and we may assume without loss of generality that $\mathbb{P}^s$ is maximal. If $\mathbb{P}^s=
\mathbb{P}({V'}_{k'}^\bot/V'_{k'})$ for some $V'_{k'}\in X'$, then in view of diagram (\ref{phi Y^2}) $\phi$ is the
constant map $V_k\mapsto V'_{k'}$, contrary to the linearity of $\phi$. Hence, $\mathbb{P}^s=
\mathbb{P}((V'_{k'+2})^*)$ for some isotropic subspace $V'_{k'+2}$ of $V'$. On the other hand,
diagram(\ref{phi Y^2}) implies that the projective space
$\mathbb{P}(V_k^*)=p_1(p_2^{-1}(V_k))$ is mapped via $\phi_{\alpha}$ to the projective space
$\mathbb{P}({V'_{k'}}^\bot/V'_{k'})=\bar{p}_1({\bar{p}_2}^{-1}(V'_{k'}))$ for $V'_{k'}=\phi(V_k)$. Thus,
$\phi_{\alpha}(\mathbb{P}(V_k^*))\subset\mathbb{P}((V'_{k'+2})^*)\cap\mathbb{P}({V'_{k'}}^\bot/V'_{k'})$. By Lemma
\ref{intersect dim=1},(ii) this implies $k=2$.  Hence, $X=GS(2,V),\ Y_{\beta}=\mathbb{P}(V)$ and, since $\phi_{\alpha}$
is linear, it is an embedding $\mathbb{P}(V)\hookrightarrow \mathbb{P}((V'_{k'+2})^*)$ corresponding to
a monomorphism $j:V\hookrightarrow(V'_{k'+2})^*$. The above mentioned reconstruction of $\phi$ from $\phi_{\alpha}$
shows now that $\phi$ decomposes as
$$
X=GS(2,V)\overset{t}\hookrightarrow G(2,V)\overset{\tilde{j}}\hookrightarrow
G(2,(V'_{k'+2})^*)\simeq G(k',V'_{k'+2})\overset{\tau}\hookrightarrow GS(k',V')=X',
$$
where $t$ is the tautological
embedding, $\tilde{j}$ is the standard extension corresponding to the monomorphism $j$, and $\tau$ is an isotropic
extension corresponding to an embedding of an isotropic subspace $V'_{k'+2}$ in $V'$. This means that $\phi$ is
a combination of isotropic and standard extensions, i.e. statement (ii) holds.

To complete case (a) it remains to consider the possibility that $\phi_{\alpha}$ is a constant map,
i.e. $\phi_{\alpha}(Y_{\beta})=\{V'_{k'+1}\}$ for some $V'_{k'+1}\subset V'$. Then diagram (\ref{phi Y^2})
implies that $\phi(X)$ lies in the projective space $\mathbb{P}(({V'}_{k'+1})^*)$ on $X'$, i.e. (iii) holds.


We now proceed to the case (b). In this case,
by Lemma \ref{intersect dim=1},(i) each space of the family $\tilde{p}:\tilde{Z}\to Y_{\beta}$ lies in a unique
projective space $\mathbb{P}(V_{k'-1}^\bot/V_{k'-1})$ of the family with base $Y'_{\beta}$. This
means that $p_{\beta}:\Sigma_{\beta}\to Y_{\beta}$ is a bijective morphism, hence an isomorphism as $Y_{\beta}$ is
a smooth variety. Therefore, there is a well-defined morphism
\begin{equation}\label{1mor i}
\phi_{\beta}:=q_{\beta}\circ p_{\beta}^{-1}:\ Y_{\beta}=GS(k-1,V)\to Y'_{\beta}=GS(k'-1,V'),
\end{equation}
and a commutative diagram similar to (\ref{phi Y^2})
\begin{equation}\label{phi Y^1}
\xymatrix{
& \Gamma\ar[rrr]^{\phi_{\Gamma}}\ar[dl]_{p_1}\ar[dr]^{p_2}&&& \Gamma'_{\beta}\ar[dl]_{p'_1}\ar[dr]^{p'_2}&
\\
Y_{\beta}\ar[dr]&& X\ar[dr]& Y'_{\beta} && X',\\
&\ar[r]_{\phi_{\beta}}&\ar[ur]&\ar[r]_{\phi}&\ar[ur]& }
\end{equation}
where $\phi,\ \Gamma,\ p_1,\ p_2,$ are as in (\ref{phi Y^2}),
$\Gamma'_{\beta}$ is the variety of isotropic $(k'-1,k')$-flags in $V'$,
and $\phi_\Gamma,\ p'_1$ and $p'_2$ are the induced projections.


Assume that $\phi_{\beta}$ is a non-constant morphism.
Then $\phi_{\beta}$ is linear, and the proof is similar to that of the linearity of $\phi_{\alpha}$.
Indeed, consider the diagram analogous to (\ref{curves}) with $\phi_{\alpha},\bar{p}_1,\bar{p}_2$ replaced respectively
by $\phi_{\beta},p'_1,p'_2$. By essentially the same argument as above, this is a commutative diagram of isomorphisms.
The fact that $\phi_{\beta}(\mathbb{P}^1_{Y_{\beta}})$ is a projective line on $Y'_{\beta}$ follows from Lemma
\ref{family CS},(i) for the data $\mathbb{P}^1=\phi(\mathbb{P}^1_X)$,
$x=\phi(\Span(\underset{V_{k-1}\in\mathbb{P}^1}\cup V_{k-1}))$, $C=\phi_{\beta}(\mathbb{P}^1_{Y_{\beta}})$.

The morphism $\phi_{\beta}$ maps a projective space $\mathbb{P}(V_k^*)$ to a unique projective space,
and thus reconstructs $\phi$ in an obvious way.


Now, by the induction assumption, (b.1) $\phi_{\beta}$ is a standard extension, or (b.2) $\phi_{\beta}$ is a
combination of isotropic and standard extensions, or (b.3) $\phi_{\beta}$ factors through a linear morphism into some
projective space $\mathbb{P}^s$ in $Y'_{\beta}$. Consider these three cases (b.1)-(b.3).

(b.1) In this case $\phi_{\beta}$ is a standard extension.
Using the reconstruction of $\phi$ via $\phi_{\beta}$ mentioned above, one immediately sees that $\phi$ is also a
standard extension.

(b.2) In this case $\phi_{\beta}$ is a combination of isotropic and standard extensions, and, using the reconstruction
of $\phi$ via $\phi_{\beta}$, the reader will check that $\phi$ also is a combination of isotropic and standard
extensions.


(b.3) In this case $\phi_{\beta}$ factors through a linear morphism of $Y_{\beta}$ into some maximal projective space
$\mathbb{P}^s$ on $Y'_\beta$.
Then $\mathbb{P}^s=\mathbb{P}^s_\beta:=\mathbb{P}({V'}_{k'-2}^\bot/V'_{k'-2})$ for some $V'_{k'-2}\subset V'$, or
$\mathbb{P}^s=\mathbb{P}^s_{\alpha}:=G(k'-1,V'_{k'})$ for some isotropic subspace $V'_{k'}\subset V'$.
The second case is clearly impossible because it would imply that
$\phi$ maps $X$ into the single point $V'_{k'}$, contrary to linearity of $\phi$. Hence,
$\mathbb{P}^s=\mathbb{P}^s_{\beta}$.

Fix $V_k\in X$ and set $V'_{k'}:=\phi(V_k).$ Diagram (\ref{phi Y^1}) shows that the projective
space $\mathbb{P}(V_k^*)=p_1(p_2^{-1}(V_k))$ is embedded by $\phi_{\beta}$ into the
intersection of the maximal projective spaces $\mathbb{P}^s_{\beta}$ and
$\mathbb{P}^{k'-1}_{\alpha}:=\mathbb{P}((V'_{k'-1})^*)=p'_1({p'}_2^{-1}(V'_{k'}))$ in $Y'_{\beta}$.
By Lemma \ref{intersect dim=1},(ii) this implies $k=2$,
i.e. $X=GS(2,V),\ Y_{\beta}=\mathbb{P}(V)$, and
$\phi_{\beta}:\ \mathbb{P}(V)\to\mathbb{P}^s_{\beta}=\mathbb{P}({V'}_{k'-2}^\bot/V'_{k'-2})$ is a linear embedding
induced by a certain monomorphism $f:V\to{V'}_{k'-2}^\bot/V'_{k'-2}$.
Diagram (\ref{phi Y^1}) shows now that $\phi$ is the composition
$$
X=GS(2,V)\overset{i}\to GS(2,{V'}_{k'-2}^\bot/V'_{k'-2})\overset{\tilde{\phi}}\hookrightarrow GS(k',V')=X',
$$
where $i$ is induced by $f$ and $\tilde{\phi}$ is the standard extension corresponding to the flag
${V'}_{k'-2}\subset{V'}_{k'-2}^\bot$ in $V'$. Being a composition of standard extensions, $\phi$ is itself a standard
extension, i.e. (i) holds.

To complete the proof in the symplectic case it remains to consider the possibility the $\phi_{\beta}$
is a constant morphism. Let $\phi_{\beta}(Y_{\beta})=\{V'_{k'-1}\}$ for some $V'_{k'-1}\subset V'$.
Then $\phi(X)$ lies in the projective space $\mathbb{P}({V'}_{k'-1}^\bot/V'_{k'-1})$ on $X'$, i.e. (iii) holds.

\vspace{2mm}
We now briefly outline the changes needed in the proof for the orthogonal case. The main idea is to replace the family
of projective spaces $PS_{\beta}(k,V)$ by the family of maximal quadrics $QO_{\beta}(k,V)$ on $X$. Note first that the
image of a quadric $Q^{\dim V-2k}_{\beta}$ under a linear morphism lies in a maximal quadric or in a projective space.
Using this and the additional conditions imposed on $k,k',\dim V,\dim V'$
we show that $\phi$ induces a well defined linear morphism of the form
\begin{equation}\label{phi a}
\phi_{\alpha}: QO_{\beta}(k,V)=GO(k-1,V)\to GO(k'+1,V')
\end{equation}
or
\begin{equation}\label{phi b}
\phi_{\beta}: QO_{\beta}(k,V)=GO(k-1,V)\to GO(k'-1,V').
\end{equation}
The above conditions ensure that $\phi$ does not map maximal quadrics of the form $Q^{\dim V-2k}_{\beta}$ into maximal
quadrics of the form $Q^4_{\gamma}$.

The linearity of $\phi_{\alpha}$ and $\phi_{\beta}$, provided that they are
non-constant morphisms, is proved by arguments similar to the above using Lemma \ref{family CO} instead of Lemma
\ref{family CS}.
The rest of the proof goes along the same lines as in the symplectic case. When working with maximal quadrics
$Q^{\dim V-2k}_{\beta}$ on $GO(k,V)$ instead of maximal projective spaces $\mathbb{P}^{2n-2k+1}_{\beta}$ on $GS(k,V)$,
one uses Lemmas \ref{Bkn},(iv) and \ref{two possib},(ii) instead of Lemmas \ref{Bkn sympl},(ii) and
\ref{intersect dim=1},(ii).

\vspace{2mm}
Finally, we leave the case $X\simeq G(k,V)$ and $X'\simeq G(k',V')$ entirely to the reader.
\end{proof}

\begin{corollary}\label{new 3.8}
\textit{Let $X\simeq G(k,V)$, $X'\simeq G(k',V')$, or $X=GO(k,V),\ X'=GO(k',V')$, or
$X=GS(k,V),\ X'=GS(k',V')$, and let $\phi: X \to X'$ be a linear morphism.
If $X=GO(k,V),\ X'=GO(k',V')$, assume in addition that either $k\le[\frac{\dim V}{2}]-3$ and
$k'\le[\frac{\dim V'}{2}]-3$, or that $[\frac{\dim V'}{2}]-k'\le[\frac{\dim V}{2}]-k\le2$ and both $\dim V$ and
$\dim V'$ are odd. Then $\phi$ is an embedding unless it factors through a projective space on $X'$
or through a maximal quadric when $X'=GO(k',V')$.}
\end{corollary}




\begin{corollary}\label{3.8}
\textit{Let $X\simeq G(k,V)$, $X'\simeq G(k',V')$, or $X=GO(k,V),\ X'=GO(k',h')$, or
$X=GS(k,V),\ X'=GS(k',V')$, and let $\phi: X \to X'$ be a linear embedding.
If $X=GO(k,V),\ X'=GO(k',V')$, assume in addition that either $k\le[\frac{\dim V}{2}]-3$ and
$k'\le[\frac{\dim V'}{2}]-3$, or that $[\frac{\dim V'}{2}]-k'\le[\frac{\dim V}{2}]-k\le2$ and both $\dim V$ and
$\dim V'$ are odd. Then some of the following statements holds:\\
(i) $\phi$ is a standard extension;\\
(ii) $X$ and $X'$ are isotropic grassmannians and $\phi$ is a combination of isotropic and standard
extensions;\\
(iii) $\phi$ factors through a projective space on $X'$ or, in case $X'=GO(k',V')$, through a maximal quadric
$Q_{\beta}^{\dim V'-2k'}$.}
\end{corollary}



\begin{remark}\label{X,X'}
Note that if $X\simeq G(k,V)$, $X'\simeq G(k',V')$ and $\phi:X\to X'$ is an embedding, the
statement of
Corollary \ref{3.8} simplifies as follows: $\phi$ is either a standard extension, or factors through a projective
space on $X'$ (cf. Proposition 3.1 in \cite{PT1}).
\end{remark}




\begin{remark}\label{X'}
If $X=GS(k,V)$, $X'=GS(\frac{\dim V'}{2},V')$ and $\phi:X\to X'$ is a linear morphism, then $k=\frac{\dim V}{2}$.
This follows easily from Lemmas \ref{Bkn sympl},(iii) and \ref{finite}.
\end{remark}


We will also need the following partial extension of Theorem 1.

\begin{proposition}\label{case n-2 etc}
Let $\dim V=2n\ge10$, $\dim V'=2n'$ and $\phi:X=GO(n-2,V)\to X'=GO(n'-2,V')$ be a linear embedding.
Then some of the following statements holds:\\
(i) $\phi$ is a standard extension;\\
(ii) $X$ and $X'$ are isotropic grassmannians and $\phi$ is a combination of isotropic and standard
extensions;\\
(iii) $\phi$ factors through a projective space on $X'$ or
through the grassmannian $G(n'-2,V'_{n'})\subset X'$ for a maximal isotropic subspace $V'_{n'}$ of $V'$.
\end{proposition}
\begin{proof}
Considering the image of the family $QO_{\beta}(n-2,V)$ under $\phi$, we see similarly to the proof of Theorem 1,
that at least one of the following morphisms
\begin{equation}\label{phi alpha}
\phi_{\alpha}: QO_{\beta}(n-2,V)=GO(n-3,V)\to PO_{\alpha}(n'-2,V'),
\end{equation}
\begin{equation}\label{phi beta}
\phi_{\beta}: QO_{\beta}(n-2,V)=GO(n-3,V)\to GO(n'-3,V'),
\end{equation}
\begin{equation}\label{phi gamma}
\phi_{\gamma}: QO_{\beta}(n-2,V)=GO(n-3,V)\to QO_{\gamma}(n'-2,V')
\end{equation}
must be well defined.

Assume that $\phi_{\alpha}$ is well defined. Then one sees that an obvious analog of diagram (\ref{phi Y^2})
applies also in the case we consider here. Set $V'_{n'-2}:=\phi(V_{n-2})$ for $V_{n-2}\in X$.
Note that $p_2^{-1}(V_{n-2})=\mathbb{P}(V_{n-2}^*)$is mapped under $\phi_{\Gamma}$ into
${p'}_2^{-1}(V'_{n'-2})\simeq\mathbb{P}^1\times\mathbb{P}^1$. Since $n\ge5$,
this map is a constant map. Hence $\phi_{\alpha}$ maps the projective space $\mathbb{P}(V_{n-2}^*)$ into a point.
Lemma \ref{finite} implies now that $\phi_{\alpha}$ is a constant map.
i.e. $\phi_{\alpha}(QO_{\beta}(n-2,V))=\{V'_{n'-1}\}$ for some $V'_{n'-1}\subset V'$. Then the analog of diagram
(\ref{phi Y^2}) implies that $\phi(X)$ lies in the projective space $\mathbb{P}((V'_{n'-1})^*)$ on $X'$, i.e.
statement (iii) holds.

Next, if $\phi_{\beta}$ is well defined, then one applies Theorem 1 to $\phi_{\beta}$ and recovers $\phi$ from
$\phi_{\beta}$ as in the proof of Theorem 1.


In the remainder of the proof we assume that $\phi_{\gamma}$ is well defined. We start by constructing a diagram
analogous to (\ref{phi Y^2}):
\begin{equation}\label{phi Y^3}
\xymatrix{
& \bar{\Gamma}\ar[rrr]^{\phi_{\bar{\Gamma}}}\ar[dl]_{\pi_1}\ar[dr]^{\pi_2} &&&
\bar{\Gamma}'\ar[dl]_{\pi'_1}\ar[dr]^{\pi'_2}& \\
Y\ar[dr]&& X\ar[dr]& Y' && X'.\\
&\ar[r]_{\phi_Y}&\ar[ur]&\ar[r]_{\phi}&\ar[ur]& }
\end{equation}
By definition, $\bar{\Gamma}$ is a fixed connected component of the variety of isotropic $(n-2,n)$-flags in $V$, and
$\bar{\Gamma}'$ is a fixed connected component of the variety of isotropic $(n'-2,n')$-flags in $V$. Next, we define
$Y$. For this we fix codimension 1 subspace $\tilde{V}$ in $V$ such that the symmetric form $\Phi|_{\tilde{V}}$ is
non-degenerate, and set $Y:=GO(n-1,\tilde{V})$. Similarly we define $Y'$ as $GO(n'-1,\tilde{V}')$. The projections
$\pi_1$, $\pi_2$, $\pi'_1$, $\pi'_2$ are as follows:
$\pi_1:(V_{n-2}\subset V_n)\mapsto V_n\cap\tilde{V}$, $\pi_2:(V_{n-2}\subset V_n)\mapsto V_{n-2}$,
$\pi'_1:(V'_{n'-2}\subset V'_{n'})\mapsto V'_{n'}\cap\tilde{V}'$, $\pi'_2:(V'_{n'-2}\subset V'_{n'})\mapsto V'_{n'-2}$.
To define the morphisms  $\phi_Y$ and $\phi_{\bar{\Gamma}}$, consider
a point $V_n\cap\tilde{V}\in Y$. By construction, the fibre $\pi_1^{-1}(V_n\cap\tilde{V})$ is isomorphic to the
grassmannian $G(n-2,V_n)$ which is isomorphically mapped onto $\pi_2(G(n-2,V_n))$. The composition
$G(n-2,V_n)\overset{\pi_2}\to \pi_2(G(n-2,V_n))\overset{\phi}\hookrightarrow X'
\overset{t} \hookrightarrow G(n'-2,V')$, where $t$ is the tautological embedding, is a linear embedding of
grassmannians, hence by Theorem 1 it is either a standard extension or factors through an embedding into a projective
space. In both cases one sees that there is a unique isotropic subspace $V'_{n'}$ of $V'$ such that
$(\phi\circ\pi_2)(G(n-2,V_n))\subset G(n'-2,V'_{n'})$. Define now $\phi_Y:Y\to Y'$ by setting
$\phi_Y(V_n\cap\tilde{V})=V'_{n'}\cap\tilde{V}'$. The morphism $\phi_{\bar{\Gamma}}:\bar{\Gamma}\to\bar{\Gamma}'$ is
then recovered by the commutativity of diagram (\ref{phi Y^3}).

Assume now that the morphism $\phi_Y$ is finite. Consider a point $V_{n-2}\in X$ and set $V'_{n'-2}=\phi(V_{n-2})$.
By diagram (\ref{phi Y^3}) the projective line $\mathbb{P}^1:=\pi_1(\pi_2^{-1}(V_{n-2}))$ on $Y$ is mapped into the
projective line $\mathbb{P}'^1:=\pi'_1({\pi'}_2^{-1}(V'_{n'-2}))$ on $Y'$. Since the morphism $\phi_Y|_{\mathbb{P}^1}$
is finite, it follows that this morphism is surjective. This implies that the morphism
$\phi_{\bar{\Gamma}}:\bar{\Gamma}\to\bar{\Gamma}'$ maps fibres of  $\pi_2$ onto fibres of $\pi'_2$.

Next, fix a point $V_{n-3}\in GO(n-3,V)$. The maximal quadric $GO(1,V_{n-3}^\bot/V_{n-3})$ is mapped by $\phi$ onto the
quadric $Q^4_{\gamma}$ corresponding to the isotropic flag $\phi_{\gamma}(V_{n-3})$. Consequently, according to
the above stated property of $\phi_{\bar{\Gamma}}$ the variety $\pi_2^{-1}(GO(1,V_{n-3}^\bot/V_{n-3}))$ is mapped by
$\phi_{\bar{\Gamma}}$ onto the variety ${\pi'}_2^{-1}(Q^4_{\gamma})$. Hence
$\pi_1(\pi_2^{-1}(GO(1,V_{n-3}^\bot/V_{n-3})))$ is mapped by $\phi_Y$ onto $\pi'_1({\pi'}_2^{-1}(Q^4_{\gamma}))$.
However, one can check that the variety $\pi_1(\pi_2^{-1}(GO(1,V_{n-3}^\bot/V_{n-3})))$ is isomorphic to $\mathbb{P}^3$,
while the variety $\pi'_1({\pi'}_2^{-1}(Q^4_{\gamma}))$ is 5-dimensional. This is a contradiction.

Hence $\phi_Y$ is not finite, and Lemma \ref{finite} implies that $\phi_Y$ is a constant map. Set $V'_{n'}=\phi_Y(Y)$.
Then diagram (\ref{phi Y^3}) yields that $\phi(X)\subset \pi'_2({\pi'}_1^{-1}(V'_{n'-2}))=G(n'-2,V'_{n'})$, and
statement (iii) holds.
\end{proof}











\vspace{1cm}

\section{Linear ind-grassmannians}\label{stand Gr}

\vspace{1cm}


\noindent



Recall that an \textit{ind-variety} is the direct limit
$\mathbf{X}=\underset{\longrightarrow}\lim X_m$
of a chain of morphisms of algebraic varieties
\begin{equation}\label{eq1}
X_1\stackrel{\phi_1}{\to}X_2\stackrel{\phi_2}{\to}\dots\stackrel{\phi_{m-1}}{\to}X_m\stackrel{\phi_m}{\to}
X_{m+1}\stackrel{\phi_{m+1}}{\to}\dots\ .
\end{equation}
Note that the direct limit of the chain (\ref{eq1})
does not change if we replace the sequence
$\{X_m\}_{m\ge1}$ by a subsequence $\{X_{i_m}\}_{m\ge1}$ and the morphisms $\phi_m$ by the compositions
$\tilde{\phi}_{i_m}:={\phi}_{i_{m+1}-1}\circ...\circ{\phi}_{i_m+1}\circ{\phi}_{i_m}$.
Let $\mathbf{X}$ be the direct limit of (\ref{eq1}) and  $\mathbf{X}'$ be the direct limit of a chain
\begin{equation}\label{eq1'}
X'_1\stackrel{\phi'_1}{\to}X'_2\stackrel{\phi'_2}{\to}\dots\stackrel{\phi'_{m-1}}{\to}X'_m\stackrel{\phi'_m}{\to}
X'_{m+1}\stackrel{\phi'_{m+1}}{\to}\dots\ .
\end{equation}
A {\it morphism of ind-varieties} $\mathbf{f}:\mathbf{X}\to\mathbf{X}'$ is a map from
$\underset{\to}\lim X_n$ to $\underset{\to}\lim X'_n$ induced by a collection of morphisms of algebraic varieties
$\{f_m:X_m\to Y_{n_m}\}_{m\ge1}$
such that $\psi_{n_m}\circ f_m=f_{m+1}\circ\phi_m$ for all $m\ge1$.
The identity morphism $\mathrm{id}_\mathbf{X}$ is a morphism which induces the identity as a set-theoretic map from
$\mathbf{X}$ to $\mathbf{X}'$.
A morphism $\mathbf{f}:\mathbf{X}\to\mathbf{X}'$ is an \textit{isomorphism} if there exists a morphism
$\mathbf{g}:\mathbf{X}'\to\mathbf{X}$ such that $\mathbf{g}\circ\mathbf{f}=\mathrm{id}_\mathbf{X}$ and
$\mathbf{f}\circ\mathbf{g}=\mathrm{id}_{\mathbf{X}'}$.

In what follows we only consider chains (\ref{eq1}) such that $X_m$ are complete algebraic varieties,
$\underset{n\to\infty}\lim(\dim X_n)=\infty$, and the morphisms $\phi_m$ are embeddings. We call such ind-varieties
\textit{locally complete}. Furthermore, we call a morphism
$\mathbf{f}:\mathbf{X}=\underset{\to}\lim X_n\to\mathbf{X}'=\underset{\to}\lim X'_n$ of locally complete
ind-varieties an \textit{embedding} if all morphisms $f_m:X_m\to X'_{n_m},\ m\ge1,$ are embeddings.
\begin{definition}\label{lin ind-Gr}
A \textit{linear ind-grassmannian} is an ind-variety $\mathbf{X}$ obtained as a direct limit of a chain of embeddings
$$
X_1\overset{\phi_1}{\hookrightarrow}X_2\overset{\phi_2}{\hookrightarrow}\dots\overset{\phi_{m-1}}{\hookrightarrow}
X_m\overset{\phi_m}{\hookrightarrow}X_{m+1}\overset{\phi_{m+1}}{\hookrightarrow}\dots
$$
where each $X_m$ is a grassmannian or an isotropic grassmannian, $\underset{n\to\infty}\lim(\dim X_n)=\infty$,
and all embeddings $\phi_m$ are linear morphisms.
\end{definition}
Note that Definition \ref{lin ind-Gr} allows for a "mixture" of all three types of grassmannians (usual grassmannians,
orthogonal grassmannians, symplectic grassmannians). Note also that when considering orthogonal grassmannians we
restrict ourselves to connected orthogonal grassmannians with Picard group isomorphic to $\mathbb{Z}$, see
\ref{linear ortho ind-grassm}.


We now define certain standard grassmannians and isotropic grassmannians.

\begin{definition}\label{G(k)}
Fix an infinite chain of vector spaces
$$
V_{n_1}\subset V_{n_2}\subset...\subset V_{n_m}\subset V_{n_{m+1}}\subset...
$$
of dimensions $n_m,\ n_m< n_{m+1}$.

a) For an integer $k$, $1\le k< n_1$, set $\mathbf{G}(k):=\underset{\to}\lim G(k,V_{n_m})$ where
$$
G(k,V_{n_1})\hookrightarrow G(k,V_{n_2})\hookrightarrow...\hookrightarrow G(k,V_{n_m})\hookrightarrow
G(k,V_{n_{m+1}})\hookrightarrow...
$$
is the chain of canonical inclusions of grassmannians.

b) For a sequence of integers $1\le k_1< k_2<...$ such that $ k_m< n_m,$  $\underset{m\to\infty}\lim(n_m-k_m)=\infty$,
set
$\mathbf{G}(\infty):=\underset{\to}\lim G(k_m,V_{n_m})$ where
$$
G(k_1,V_{n_1})\hookrightarrow G(k_2,V_{n_2})\hookrightarrow...\hookrightarrow G(k_m,V_{n_m})\hookrightarrow
G(k_{m+1},V_{n_{m+1}})\hookrightarrow...
$$
is an arbitrary chain of standard extensions of grassmannians.

c) Assume that $V_{n_m}$ are endowed with compatible non-degenerate symmetric (respectively, symplectic) forms $\Phi_m$.
In the symplectic case $\frac{1}{2}n_m\in\mathbb{Z}_+$.
For an integer $k,\ 1\le k\le[\frac{n_1}{2}]$, set
 $\mathbf{G}\mathrm{O}(k,\infty):=\underset{\to}\lim GO(k,V_{n_m})$ (respectively,
 $\mathbf{G}\mathrm{S}(k,\infty):=\underset{\to}\lim GS(k,V_{n_m})$) where
$$
GO(k,V_{n_1})\hookrightarrow GO(k,V_{n_2})\hookrightarrow...\hookrightarrow GO(k,V_{n_m})\hookrightarrow
GO(k,V_{n_{m+1}})\hookrightarrow...
$$
(respectively,
$$
GS(k,V_{n_1})\hookrightarrow GS(k,V_{n_2})\hookrightarrow...\hookrightarrow GS(k,V_{n_m})\hookrightarrow
GS(k,V_{n_{m+1}})\hookrightarrow...)
$$
is the chain of canonical inclusions of isotropic grassmannians.

d) For a sequence of integers $1\le k_1< k_2<...$ such that $k_m<[\frac{n_m}{2}],$
$\underset{m\to\infty}\lim([\frac{n_m}{2}]-k_m)=\infty$,
set $\mathbf{G}\mathrm{O}(\infty,\infty)=\underset{\to}\lim G\mathrm{O}(k_m,V_{n_m})$
(respectively, $\mathbf{G}\mathrm{S}(\infty,\infty):=\underset{\to}\lim GS(k_m,V_{n_m})$)
where
\begin{equation}\label{chain O}
GO(k_1,V_{n_1})\hookrightarrow GO(k_2,V_{n_2})\hookrightarrow...\hookrightarrow GO(k_m,V_{n_m})\hookrightarrow
GO(k_{m+1},V_{n_{m+1}})\hookrightarrow...
\end{equation}
(respectively,
\begin{equation}\label{chain S}
GS(k_1,V_{n_1})\hookrightarrow GS(k_2,V_{n_2})\hookrightarrow...\hookrightarrow GS(k_m,V_{n_m})\hookrightarrow
GS(k_{m+1},V_{n_{m+1}})\hookrightarrow...)
\end{equation}
is an arbitrary chain of standard extensions of isotropic grassmannians.

e) In the symplectic case, consider a sequence of integers $1\le k_1< k_2<...$ such that $k_m<\frac{n_m}{2}$,
$\underset{m\to\infty}\lim(\frac{n_m}{2}-k_m)=k\in\mathbb{N}$, and set
$\mathbf{G}\mathrm{S}(\infty,k):=\underset{\to}\lim GS(k_m,V_{n_m})$
for any chain of standard extensions (\ref{chain S}). In the orthogonal case, assume first that $\dim V_{n_m}$ are even.
Then set $\mathbf{G}\mathrm{O}^0(\infty,k):=\underset{\to}\lim GO(k_m,V_{n_m})$ for a chain (\ref{chain O}) where
$k_m<\frac{n_m}{2}$, $\underset{m\to\infty}\lim(\frac{n_m}{2}-k_m)=k\in\mathbb{N},\ k\ge2$. Finally, consider the
orthogonal case under the assumption that $\dim V_{n_m}$ are odd. Then set
$\mathbf{G}\mathrm{O}^1(\infty,k):=\underset{\to}\lim GO(k_m,V_{n_m})$ for a chain (\ref{chain O}) where
$k_m\leq [\frac{n_m}{2}]$, $\underset{m\to\infty}\lim([\frac{n_m}{2}]-k_m)=k\in\mathbb{N}$.
\end{definition}


The \textit{infinite projective space} $\mathbf{P}^\infty$ is defined as the ind-variety $\mathbf{G}(1)$.
Note that $\mathbf{P}^\infty\simeq\mathbf{G}\mathrm{S}(1)$. When writing $\mathbf{G}\mathrm{O}^0(\infty,k)$ below
we automatically assume $k\ne1$.


\begin{lemma}\label{G infty}
All standard ind-grassmannians $\mathbf{G}(\infty)$, $\mathbf{G}\mathrm{O}(\infty,\infty)$,
$\mathbf{G}\mathrm{S}(\infty,\infty)$, $\mathbf{G}(k)$, $\mathbf{G}\mathrm{O}(k,\infty)$,
$\mathbf{G}\mathrm{S}(k,\infty)$, $\mathbf{G}\mathrm{O}^0(\infty,k)$, $\mathbf{G}\mathrm{O}^1(\infty,k)$,
$\mathbf{G}\mathrm{S}(\infty,k)$,
are well defined. In other words, a standard grassmannian does not depend,
up to an isomorphism of ind-varieties, on the specific chain of standard embeddings used in its definition.
\end{lemma}

\begin{proof}
We consider only $\mathbf{G}(\infty)$. All other cases are similar.
Let two chains of strict standard extensions
$$
G(k_1,V_{n_1})\overset{\phi_1}\hookrightarrow G(k_2,V_{n_2})\overset{\phi_2}\hookrightarrow...
\overset{\phi_{m-1}}\hookrightarrow G(k_m,V_{n_m})\overset{\phi_m}\hookrightarrow G(k_{m+1},V_{n_{m+1}}
)\overset{\phi_{m+1}}\hookrightarrow...,
$$
$$
G(k'_1,V_{n'_1})\overset{\phi'_1}\hookrightarrow G(k'_2,V_{n'_2})\overset{\phi'_2}\hookrightarrow...
\overset{\phi'_{m-1}}\hookrightarrow G(k'_m,V_{n'_m})\overset{\phi'_m}\hookrightarrow G(k'_{m+1},V_{n'_{m+1}}
)\overset{\phi'_{m+1}}\hookrightarrow...,
$$
such that
$$
\underset{m\to\infty}\lim k_m=\underset{m\to\infty}\lim k'_m=\underset{m\to\infty}\lim(n_m-k_m)=
\underset{m\to\infty}\lim(n'_m-k'_m)=\infty,
$$
be given. We will show that their respective direct limits $\mathbf{G}(\infty)$ and $\mathbf{G}'(\infty)$ are
isomorphic as ind-varieties.

For this, we have to construct two infinite subsequences
$\{i_s\}_{s\ge1}$ and $\{j_s\}_{s\ge1}$ of $\mathbb{Z}_+$
and two sets of morphisms
$\mathbf{f}=\{f_s:G(k_{i_s},V_{n_{i_s}})\to G(k'_{j_s},V'_{n'_{j_s}})\}_{s\ge1}$,
$\mathbf{g}=\{g_m:G(k'_{j_s},V'_{n'_{j_s}})\to G(k_{i_{s+1}},V_{n_{i_{s+1}}})\}_{m\ge1}$
such that they determine morphisms of ind-varieties $\mathbf{f}:\mathbf{G}(\infty)\to\mathbf{G}'(\infty),
\mathbf{g}:\mathbf{G}'(\infty)\to\mathbf{G}(\infty)$ with
$\mathbf{g}\circ\mathbf{f}=\mathrm{id}_{\mathbf{G}(\infty)}$ and
$\mathbf{f}\circ\mathbf{g}=\mathrm{id}_{\mathbf{G}'(\infty)}$.
Assume that the desired subsequences $\{i_s\}_{s\ge1},\{j_s\}_{s\ge1}$ and
morphisms $f_l,g_l$ are constructed for $1\le l\le s-1$,
and that these morphisms are strict standard extensions. Denote for short
$k:=k_{i_s},\ n:=n_{i_s},\ V:=V_n,\ k':=k'_{j_s},\ n':=n'_{j_s},\ V':=V'_{n'},\
G:=G(k,V),\ G':=G(k',V'),\ f:=f_s:G\hookrightarrow G',\ \tilde{k}:=k_{i_{s+1}},
\tilde{n}:=n_{i_{s+1}},\ \tilde{V}:=V_{\tilde{n}},\ \tilde{G}:=G(\tilde{k},\tilde{V}),\
\phi:=\phi_{i_s}:G\hookrightarrow\tilde{G}$. Without loss of generality that we assume that $\tilde{k}>k'$.
By Remark \ref{intrinsic def}, $f$ is given by a triple $(W_f,U_f,\underline{f})$, where $W_f\subset U_f$
is a flag in $V'$. Respectively, $\phi$ is given by a triple $(W_\phi,U_\phi,\underline{\phi})$,
where $W_\phi\subset U_\phi$ is a flag in $\tilde{V}$.

For the induction step we will now find a strict standard extension
$g:=g_s:G'\hookrightarrow\tilde{G}$ such that $g\circ f=\phi$. Indeed,
consider the exact triples
$0\to W_f\to U_f\xrightarrow{\underline{f}}V\to0,\ \ \
0\to W_{\phi}\to U_{\phi}\xrightarrow{\underline{\phi}}V\to0$.
Since both $\underline{f}$ and $\underline{\phi}$ are epimorphisms, and $\dim U_{\phi}>\dim U_f$ as $\tilde{k}>k'$,
it follows that there exists a (non-unique) epimorphism $\varepsilon_U:U_{\phi}\twoheadrightarrow U_f$
such that $\underline{\phi}=\underline{f}\circ\varepsilon_U$.
Then $\varepsilon_U|_W$ is a well-defined epimorphism $W_{\phi}\twoheadrightarrow W_f$.
Putting $W_g:=\ker{\varepsilon_U}$, we have the exact triple
$0\to W_g\to U_{\phi}\xrightarrow{\varepsilon_U}U_f\to0$.
Next, set
$U'_g:=W_g\oplus V'$, and fix embeddings $i:U'_g\hookrightarrow\tilde{V}$ and $j:U_\phi\hookrightarrow U'_g$ such that the composition $i\circ j$ induces the identity on $U_\phi$, and such that $\text{pr}_2\circ j = \varepsilon_U$ where $\text{pr}_2$ is the projection onto the second direct summand.   It is easy to check that such embeddings $i$ and $j$ exist.
Then $W_g\subset U_g:=i(U'_g)$ is a flag in $\tilde{V}$ equipped with an isomorphism $\underline{g}:U_g/W_g\simeq V'$.
The corresponding strict standard extension $g:G'\hookrightarrow\tilde{G}$ satisfies the property $g\circ f=\phi$,
as claimed.
\end{proof}


Note furthermore that the standard ind-grassmannians introduced above are isomorphic to certain ind-varieties
introduced in \cite{DiP}. More precisely, let $\widetilde{V}$ be a countable-dimensional vector space with basis
$\{v_1,...,v_n,...\}$ and let
$\widetilde{W}\subset\widetilde{V}$ be a subspace generated by a subset of $\{v_1,...,v_n,...\}$. Then
$G(\widetilde{W},\widetilde{V})$ is by definition the set of subspaces $\widetilde{E}\subset\widetilde{V}$
satisfying the following two conditions:

(i) $\Span(\{v_1,...,v_n,...\}\cap\widetilde{E})$ is of finite codimension in $\widetilde{E}$;

(ii) there exists a finite-dimensional subspace $\widetilde{U}\subset\widetilde{V}$ such that
$\widetilde{W}\subset\widetilde{E}+\widetilde{U}$, $\widetilde{E}\subset\widetilde{W}+\widetilde{U}$,
$\dim(\widetilde{E}\cap\widetilde{U})=\dim(\widetilde{W}\cap\widetilde{U})$.\\
Then it is easy to see (a much stronger result is proved in \cite{DiP}) that $G(\widetilde{W},\widetilde{V})$
has a natural structure of an ind-variety such that $G(\widetilde{W},\widetilde{V})$ is the direct limit of a chain of
standard extensions of grassmannians. Moreover,
$$
G(\widetilde{W},\widetilde{V})\cong\mathbf{G}(\min\{\dim\widetilde{W},\codim_{\widetilde{V}}\widetilde{W}\}).
$$
Similarly, in the isotropic case (i.e. in the case when $\widetilde{W}$ is equipped with an appropriate
non-degenerate quadratic form)
the standard isotropic ind-grassmannians introduced in this paper represent all isomorphism
classes of ind-varieties $G(\widetilde{W},\widetilde{V})$  introduced in \cite{DiP}
(in this case $\widetilde{W}$ is an isotropic subspace of $\widetilde{V}$) and satisfying
$\Pic G(\widetilde{W},\widetilde{V})\simeq\mathbb{Z}$.





\vspace{1cm}

\section{Classification of linear ind-grassmannians}\label{sec5}

\vspace{0.5cm}

In this section we prove the following main result of the note.

\textbf{Theorem 2}. \textit{Every linear ind-grassmannian is isomorphic as an ind-variety to one of the standard
ind-grassmannians $\mathbf{G}(k)$ for $k\ge1$, $\mathbf{G}(\infty)$, $\mathbf{G}\mathrm{O}(k,\infty)$ for $k\ge1$,
$\mathbf{G}\mathrm{O}^0(\infty,k)$ for $k\ge2$, $\mathbf{G}\mathrm{O}^1(\infty,k)$ for $k\ge0$,
$\mathbf{G}\mathrm{O}(\infty,\infty)$, $\mathbf{G}\mathrm{S}(k,\infty)$ for $k\ge2$, $\mathbf{G}\mathrm{S}(\infty,k)$
for $k\ge0$, $\mathbf{G}\mathrm{S}(\infty,\infty)$, and the latter are pairwise non-isomorphic}.

\begin{proof}
Let a linear ind-grassmannian $\mathbf{X}$ be given as the direct limit of a chain of embeddings
$$
X_1\stackrel{\phi_1}{\hookrightarrow}X_2\stackrel{\phi_2}{\hookrightarrow}
\dots\stackrel{\phi_{m-1}}{\hookrightarrow}X_m\stackrel{\phi_m}{\hookrightarrow}
X_{m+1}\stackrel{\phi_{m+1}}{\hookrightarrow}\dots,
$$
where $X_m$ are grassmannians, possibly orthogonal or symplectic, such that
$\underset{m\to\infty}\lim(\dim X_m)=\infty.$
Then, for infinitely many $m$, $X_m$ will be a grassmannian, or an orthogonal grassmannian, or a symplectic
grassmannian. Therefore, without loss of generality, we can assume that all $X_m$ are of one of the above three types.

Suppose first that all $X_m$ are grassmannians. Then we have the following two options: for infinitely many $m$,
the embedding $\phi_m:X_m\to X_{m+1}$ factors through an embedding of a projective space into $X_{m+1}$, i.e.
there exists a commutative diagram of embeddings
$$
\xymatrix{
X_m \ar[rr]^-{\phi_m}\ar[dr]& & X_{m+1}\\
& \mathbb{P}^{j_m},\ar[ur] & }
$$
or this is not the case. In the first case $\mathbf{X}\simeq\underset{\to}\lim~\mathbb{P}^{j_m}$, hence
$\mathbf{X}\simeq\mathbf{P}^{\infty}$.
In the second case, by deleting some first embeddings we can assume that none of the embeddings $\phi_m:X_m\to X_{m+1}$
factors through an embedding of a projective space into $X_{m+1}$. Then, Corollary \ref{3.8} implies that all embeddings
$\phi_m$ are standard extensions, hence $\mathbf{X}$ is isomorphic to $\mathbf{G}(k)$ or $\mathbf{G}(\infty)$.


In the symplectic case, the reader will argue in a similar way that Corollary \ref{3.8} implies that $\mathbf{X}$ is
either isomorphic to $\mathbf{G}(k)$ or $\mathbf{G}(\infty)$ (this happens when all $\phi_m$ are combinations of
isotropic and standard extensions or factor through projective spaces), or to one of the standard symplectic
ind-grassmannians.

The orthogonal case is similar but has some special features. First, if all morphisms $\phi_m$ factor through
respective quadrics $Q_{\beta}^{\dim V_{m+1}-2k_{m+1}}$, one needs to prove that the direct limit of any chain
of linear embeddings
$$
Q_1{\hookrightarrow}Q_2{\hookrightarrow}\dots{\hookrightarrow}Q_m{\hookrightarrow}Q_{m+1}{\hookrightarrow}\dots\ ,
$$
where $\underset{m\to\infty}\lim\dim Q_m=\infty$, is isomorphic either to $\mathbf{P}^{\infty}$ or to
$\mathbf{G}\mathrm{O}(1,\infty)$. This is an exercise which we leave to the reader. Second, in the orthogonal case one
applies Corollary \ref{3.8} when $[\frac{\dim V_m}{2}]-k_m\ge3$ for infinitely many $m$ (in this case one can assume
without loss of generality that $[\frac{\dim V_m}{2}]-k_m\ge3$ for all $m$). The case when
$[\frac{\dim V_m}{2}]-k_m\le2$ for infinitely many $m$ needs special attention. In the latter case one assumes without
loss of generality that $[\frac{\dim V_m}{2}]-k_m$ is constant and then applies Theorem 1 when $\dim V_m$ is odd for
all $m$, and Proposition \ref{case n-2 etc} when $\dim V_m$ is even for all $m$ (in the latter case
$\frac{\dim V_m}{2}-k_m=2$ for all $m$).

The first claim of Theorem 2 is now proved.

The claim that the standard ind-grassmannians are pairwise non-isomorphic follows from Lemmas \ref{non-isom1},
\ref{non-isom2} and \ref{non-isom3} below.
\end{proof}

In what follows we will sometimes write $\mathbf{G}\mathrm{O}(\infty,k)$ meaning $\mathbf{G}\mathrm{O}^0(\infty,k)$
or $\mathbf{G}\mathrm{O}^1(\infty,k)$. This allows the simultaneous consideration of $\mathbf{G}\mathrm{O}^0(\infty,k)$
and $\mathbf{G}\mathrm{O}^1(\infty,k)$.

\begin{lemma}\label{non-isom1}
(i) Let $k,k'\in\mathbb{Z}_{+}\cup\{\infty\},\ k\ne k'$. Then
$\mathbf{G}(k)\not\simeq\mathbf{G}(k')$, $\mathbf{G}\mathrm{O}(k,\infty)\not\simeq\mathbf{G}\mathrm{O}(k',\infty)$,
$\mathbf{G}\mathrm{O}(k,\infty)\not\simeq\mathbf{G}\mathrm{O}(k',\infty)$.

(ii) Let $k\ge2$. Then $\mathbf{G}\mathrm{O}^0(\infty,k)\not\simeq\mathbf{G}\mathrm{O}^1(\infty,k)$.

(iii) Let $k,k'\in\mathbb{N}\cup\{\infty\},\ k\ne k'$. Then $\mathbf{G}\mathrm{O}(\infty,k)\not\simeq
\mathbf{G}\mathrm{O}(\infty,k')$, $\mathbf{G}\mathrm{S}(\infty,k)\not\simeq\mathbf{G}\mathrm{S}(\infty,k')$.

(iv) Let $k\in\mathbb{N}\cup\{\infty\},\ k'\in\mathbb{Z}_{+}\cup\{\infty\},\ k\ne k'$.
Then $\mathbf{G}\mathrm{O}(\infty,k)\not\simeq\mathbf{G}\mathrm{O}(k',\infty)$, $\mathbf{G}\mathrm{S}(\infty,k)
\not\simeq\mathbf{G}\mathrm{S}(k',\infty)$.

\end{lemma}
\begin{proof}
In (i), (iii) and (iv) we only consider the symplectic case and leave the other cases to the reader.

(i) Let $k>k'$. Assume that $k\in\mathbb{Z}_{+}$ and that $\mathbf{X}:=\mathbf{G}\mathrm{S}(k,\infty)$ and
$\mathbf{X}':=\mathbf{G}\mathrm{S}(k',\infty)$ are isomorphic. This implies that there exist
subsequences $\{i_s\}_{s\ge1}$ and $\{j_s\}_{s\ge1}$ of $\mathbb{Z}_{+}$ and
a chain of linear embeddings
\begin{equation}\label{chain1}
...\hookrightarrow GS(k,V_{n_{i_s}})\overset{f_s}\hookrightarrow GS(k',V'_{n'_{j_s}})\overset{g_s}\hookrightarrow
GS(k,V_{n_{i_{s+1}}})\overset{f_{s+1}}\hookrightarrow GS(k',V'_{n'_{j_{s+1}}})\overset{g_{s+1}}\hookrightarrow ...\ ,
\end{equation}
such that the compositions $g_s\circ f_s$ and $f_{s+1}\circ g_s$ are standard extensions and
the direct limit of the chain (\ref{chain1}) is isomorphic to both $\mathbf{X}$ and $\mathbf{X}'$.
According to Corollary \ref{3.8}, we can assume without loss of generality that all embeddings $f_s$ and $g_s$ are
standard extensions, or factor through isotropic extensions, or factor through embeddings to projective spaces.

In the first case, since
$GS(k,V_{n_{i_s}})\overset{f_s}\hookrightarrow GS(k',V'_{n'_{j_s}})$
is a standard extension, it follows from (\ref{stand ineq1}) that $k'\ge k$, contrary to the assumption.

In the third case both $\mathbf{X}$ and $\mathbf{X}'$ are isomorphic to $\mathbf{P}^\infty$.
On the other hand, Remark \ref{prop of st ext} implies that $\mathbf{X}$ is not isomorphic to
$\mathbf{P}^\infty$ as $k>1$.


Consider now the second case. Here $f_s$ factorizes as
$f_s: GS(k,V_{n_{i_s}})\overset{t}\hookrightarrow G(k,V_{n_{i_s}})\overset{\tilde{f}_s}\hookrightarrow
GS(k',V'_{n'_{j_s}})$,
where $t$ is the tautological embedding and $\tilde{f}_s$ is an isotropic extension followed by a standard extension.
The composition
\begin{equation}\label{GkV}
G(k,V_{n_{i_s}})\overset{\tilde{f}_s}\hookrightarrow GS(k',V'_{n'_{j_s}})\overset{\tilde{t}}\hookrightarrow
G(k',V'_{n'_{j_s}}),
\end{equation}
$\tilde{t}$ being the tautological embedding, is a standard extension or factors through a projective space.
The latter assumption leads to the same contradiction as in the above considered third case, so we must assume that
(\ref{GkV}) is a standard extension. The existence of a standard extension
$G(k,V_{n_{i_s}})\hookrightarrow G(k',V'_{n'_{j_s}})$ implies $k'\ge k,\ k'-k\le n'_{j_s}-n_{i_s},$ or
$k\le n'_{j_s}-k'$, $k+k'\ge n_{i_s}$ (see Remark \ref{prop of st ext1}). Since for $n_{i_s}$ large enough, both pairs
of inequalities contradict our assumption that $k>k'$, we conclude that the second case is also impossible.

We have now shown that all three cases lead to contradictions, hence (i) follows for $k\in\mathbb{Z}_{+}$.
The argument for $k=\infty$ is very similar.

(ii) The maximal quadrics on $\mathbf{G}\mathrm{O}^0(\infty,k)$ not lying in projective spaces have dimension $2k$,
while the maximal quadrics on $\mathbf{G}\mathrm{O}^1(\infty,k)$ not lying in projective spaces have dimension $2k+1$, see Lemma \ref{Bkn}.
This imlies that $\mathbf{G}\mathrm{O}^0(\infty,k)\not\simeq\mathbf{G}\mathrm{O}^1(\infty,k)$.

(iii) Let $\mathbb{Z}_{+}\ni k>k'$. Assume that $\mathbf{X}:=\mathbf{G}\mathrm{S}(k,\infty)$ and
$\mathbf{X}':=\mathbf{G}\mathrm{S}(k',\infty)$ are
isomorphic. As above, this implies that there exists a chain of linear embeddings (\ref{chain1})
such that the compositions $g_s\circ f_s$ and $f_{s+1}\circ g_s$ are standard extensions and
the direct limit of the chain (\ref{chain1}) is isomorphic to both $\mathbf{X}$ and $\mathbf{X}'$.
Without loss of generality we can assume that all embeddings
$f_s$ and $g_s$ are standard extensions, or factor through isotropic extensions, or factor through embeddings
to projective spaces.

In the first case we have a standard extension
$$
GS(\frac{1}{2}\dim V_{n_{i_s}}-k,V_{n_{i_s}})\overset{f_s}\hookrightarrow GS(\frac{1}{2}\dim V'_{n'_{j_s}}-
k',V'_{n'_{j_s}}),
$$
and (\ref{stand ineq2}) gives $k\le k'$, contrary to the assumption.

The arguments in the second and third case are similar to the respective arguments in (i).

The proof is finished for $k<\infty$. The case $k=\infty$ is similar.

(iv) The argument is practically the same as in (i).

\end{proof}

\begin{lemma}\label{non-isom2}
For any $k,k'\in\mathbb{Z}_{+}\cup\{\infty\},k''\in\mathbb{N}\cup\{\infty\}$ the following assertions hold.

(i) $\mathbf{G}(k)\not\simeq\mathbf{G}\mathrm{S}(k',\infty)$, unless $k=k'=1$,\
$\mathbf{G}(k)\not\simeq\mathbf{G}\mathrm{O}(k',\infty)$,

(ii) $\mathbf{G}(k)\not\simeq\mathbf{G}\mathrm{O}(\infty,k'')$,\
$\mathbf{G}(k)\not\simeq\mathbf{G}\mathrm{S}(\infty,k'')$.

\end{lemma}
\begin{proof} Again we consider only the symplectic case and leave the orthogonal case to the reader.

(i) We have to prove that $\mathbf{G}(k)\not\simeq\mathbf{G}\mathrm{S}(k',\infty)$, unless $k=k'=1$.
The case $k'=1,\ k>k'$, is already considered in Lemma \ref{non-isom1},(i), so we can assume $k\ne1,\ k'\ne1,\ k\ne k'$.

Let $\mathbf{G}(k)$ (respectively, $\mathbf{G}\mathrm{S}(k',\infty)$) be given as the direct limit of a chain of strict
standard extensions
$$
G(k,V_{n_1})\hookrightarrow G(k,V_{n_2})\hookrightarrow...\hookrightarrow G(k,V_{n_m})\hookrightarrow
G(k,V_{n_{m+1}})\hookrightarrow...
$$
(respectively,
$$
GS(k',V'_{n'_1})\overset{\phi_1}\hookrightarrow GS(k',V'_{n'_2})\overset{\phi_2}\hookrightarrow...
\overset{\phi_{m-1}}\hookrightarrow GS(k',V'_{n'_m})\overset{\phi_m}\hookrightarrow
GS(k',V'_{n'_{m+1}})\overset{\phi_{m+1}}\hookrightarrow...).
$$
Suppose that $\mathbf{G}(k)\simeq\mathbf{G}\mathrm{S}(k',\infty)$. This means
that there exist two infinite subsequences
$\{i_s\}_{s\ge1}$ and $\{j_s\}_{s\ge1}$ of $\mathbb{Z}_+$
and two sets of morphisms
$\mathbf{f}=\{f_s:G(k,V_{n_{i_s}})\to GS(k',V'_{n'_{j_s}})\}_{s\ge1}$,
$\mathbf{g}=\{g_s:GS(k',V'_{n'_{j_s}})\to G(k,V_{n_{i_{s+1}}})\}_{m\ge1}$
which determine morphisms of ind-varieties $\mathbf{f}:\mathbf{G}(k)\to\mathbf{G}\mathrm{S}(k',\infty),
\mathbf{g}:\mathbf{G}\mathrm{S}(k',\infty)\to\mathbf{G}(k)$ with
$\mathbf{g}\circ\mathbf{f}=\mathrm{id}_{\mathbf{G}(k)}$ and
$\mathbf{f}\circ\mathbf{g}=\mathrm{id}_{\mathbf{G}\mathrm{S}(k',\infty)}$.

Set $\tilde{V}:=V_{n_{i_{s+1}}},\ \tilde{G}:=G(k,\tilde{V}),\ V':=V'_{n'_{j_s}},
\ GS:=GS(k',V'),\ \tilde{V}':=V'_{n'_{j_{s+1}}},\ \widetilde{GS}:=GS(k',\tilde{V}'),\ g:=g_s:GS\hookrightarrow\tilde{G},
\ f:=f_{s+1}:\tilde{G}\hookrightarrow\widetilde{GS},
\ \phi:=\phi_{i_s}:GS\hookrightarrow\widetilde{GS}$. Note that $\phi$ is a standard extension and $\phi=f\circ g$ by
construction.

Consider the composition
$F:\tilde{G}\overset{f}\hookrightarrow \widetilde{GS}\overset{i}\hookrightarrow
G(k',\tilde{V}')$ where $i$ is the tautological embedding.

The morphism $F$ is a linear embedding, hence, by Corollary \ref{3.8}, we may assume
without loss of generality that

(a) $F$ is a standard extension,\\
or

(b) $F$ factors through an embedding into a projective space.

Consider these two cases.

(a) By Remark \ref{intrinsic def2},
$\phi$ is given by a triple $(W_{\phi},U_{\phi},\underline{\phi})$ where $W_{\phi}\subset U_{\phi}$ is a flag in
$\tilde{V}'$. Furthermore, without loss of generality we may assume that $F$ is given by a triple
$(W_f,U_f,\underline{F})$ for a flag $W_f\subset U_f$ in $\tilde{V}'$. Since $\phi(GS)=f\circ g(GS)\subset
F(\tilde{G})$, the following chain of inclusions holds:
$$
W_F\subset W_{\phi}\subset U_{\phi}\subset U_F\subset \tilde{V}'.
$$
Therefore we have an embedding $U_{\phi}/W_{\phi}\hookrightarrow U_F/W_{\phi}$ and a projection
$U_F/W_F\twoheadrightarrow U_F/W_{\phi}$. However, since $\phi$ is a standard extension, the fixed symplectic form
$\tilde{\Phi}'$ on $\tilde{V}'$ induces a nondegenerate form on $U_{\phi}/W_{\phi}$, while it induces the zero form on
$U_F/W_F$ as $f(\tilde{G})\subset\widetilde{GS}$.
This contradiction shows that the case (a) is impossible.

(b) By assumption, $F:\tilde{G}\overset{f}\hookrightarrow\widetilde{GS}\overset{i}\hookrightarrow G(k',\tilde{V}')$
decomposes as
$\tilde{G}\hookrightarrow\mathbb{P}^r\hookrightarrow G(k',\tilde{V}')$. Without loss of generality we assume that
$\mathbb{P}^r$ is a maximal projective space on $G(k',\tilde{V}')$, and consider the two possible cases:
$\mathbb{P}^r=\{V_{k'}\subset\tilde{V}'|V_{k'-1}\subset V_{k'}\subset V^\bot_{k'-1}\}$ and
$\mathbb{P}^r=\{V_{k'}\subset\tilde{V}'$
$|V_{k'}\subset V_{k'+1}\}$
for some fixed subspaces $V_{k'-1}$ and $V_{k'+1}$ of $\tilde{V}'$, $V_{k'-1}$ being isotropic.

In the former case any
$V_{k'}\in G(k',\tilde{V}')$ such that $V_{k'-1}\subset V_{k'}\subset V_{k'-1}^\bot$
is isotropic, i.e. $V_{k'}\in\widetilde{GS}\cap\mathbb{P}^r$. In other words,
$$
\widetilde{GS}\cap\mathbb{P}^r=\mathbb{P}(V_{k'-1}^\bot/V_{k'-1}),
$$
where the intersection is taken in $G(k',\tilde{V}')$. This means that $\phi$ factors through
a projective subspace of $\widetilde{GS}$, which contradicts Remark \ref{prop of st ext}.
Hence, the former case is impossible.

In the latter case it is easy to check that, for $n'_{j_{s+1}}=\dim\tilde{V}'>2$, the subspace
$V_{k'+1}\subset \tilde{V}'$ is necessarily isotropic. Then $\widetilde{GS}\cap\mathbb{P}^r=\mathbb{P}((V_{k'+1})^*)$,
and we are led to a contradiction as in the former case.

(ii) The proof is analogous to the proof of (i) and we leave it to the reader.
\end{proof}


\begin{lemma}\label{point, line}
Let $1\le k<n=[\dim V/2]$ and $\phi:GO(k,V)\to GO(k',V'),\ V_k\mapsto V_k\oplus W$, be a standard extension.
Let two maximal projective spaces $\mathbb{P}^k_{\alpha}$ and $\mathbb{P}^{n-k}_{\beta}$ intersect in a point.
Then there exist maximal projective spaces $\mathbb{P}^{k'}_{\alpha}$ and
$\mathbb{P}^{n'-k'}_{\beta},\ n'=[\dim V'/2]$, on
$GO(k',V')$ such that $\phi(\mathbb{P}^k_{\alpha})\subset\mathbb{P}^{k'}_{\alpha}$,
$\phi(\mathbb{P}^{n-k}_{\beta})\subset\mathbb{P}^{n'-k'}_{\beta}$, and
$\mathbb{P}^{k'}_{\alpha}\cap\mathbb{P}^{n'-k'}_{\beta}$ is a point.
\end{lemma}

\begin{proof}
The projective spaces $\mathbb{P}^k_{\alpha}$ and $\mathbb{P}^{n-k}_{\beta}$ determine a configuration
$V_{k-1},V_{k+1},V_n$ as in Lemma \ref{two possib},(iv). The subspaces $V_{k-1}\oplus W,V_{k+1}\oplus W,V_n\oplus W$
of $V'$ form the configuration which determines the desired projective spaces $\mathbb{P}^{k'}_{\alpha}$ and
$\mathbb{P}^{n'-k'}_{\beta}$.
\end{proof}



\begin{lemma}\label{non-isom3}

(i) $\mathbf{G}\mathrm{O}(k,\infty)\not\simeq\mathbf{G}\mathrm{S}(k',\infty)$ for $k,k'\in\mathbb{Z}_{+}\cup\{\infty\}$.

(ii) $\mathbf{G}\mathrm{O}(k,\infty)\not\simeq\mathbf{G}\mathrm{S}(\infty,k')$ for $k\in\mathbb{Z}_{+}\cup\{\infty\},
k'\in\mathbb{N}\cup\{\infty\}$.

(iii) $\mathbf{G}\mathrm{O}(\infty,k)\not\simeq\mathbf{G}\mathrm{S}(k',\infty)$ for
$k\in\mathbb{N}\cup\{\infty\},k'\in\mathbb{Z}_{+}\cup\{\infty\}$.

(iv) $\mathbf{G}\mathrm{O}(\infty,k)\not\simeq\mathbf{G}\mathrm{S}(\infty,k')$ for $k,k'\in\mathbb{N}\cup\{\infty\}$.
\end{lemma}

\begin{proof} We consider in detail only the case of  $\mathbf{G}\mathrm{O}(\infty,\infty)$ and
$\mathbf{G}\mathrm{S}(\infty,\infty)$.
Let $\mathbb{P}^q$ for $q\ge2$ be a projective space on $\mathbf{G}\mathrm{O}(\infty,\infty)$ (respectively,
$\mathbf{G}\mathrm{S}(\infty,\infty)$). We now explain how to label $\mathbb{P}^q$ as $\mathbb{P}^q_{\alpha}$ or
$\mathbb{P}^q_{\beta}$. Fix an arbitrary chain of standard extensions
\begin{equation}\label{exh1}
GO(k_1,V_{n_1})\hookrightarrow GO(k_2,V_{n_2})\hookrightarrow...\hookrightarrow GO(k_m,V_{n_m})\hookrightarrow
GO(k_{m+1},V_{n_{m+1}})\hookrightarrow...
\end{equation}
(respectively,
\begin{equation}\label{exh2}
GS(k'_1,V_{n'_1})\hookrightarrow GS(k'_2,V_{n'_2})\hookrightarrow...\hookrightarrow GS(k'_m,V_{n'_m})\hookrightarrow
GS(k'_{m+1},V_{n'_{m+1}})\hookrightarrow...)
\end{equation}
such that
$$
\underset{m\to\infty}\lim k_m=\underset{m\to\infty}\lim(n_m-k_m)=\infty
$$
(respectively,
$$
\underset{m\to\infty}\lim k'_m=\underset{m\to\infty}\lim(n'_m-k'_m)=\infty)
$$
and
$\underset{\to}\lim GO(k_m,V_{n_m})=\mathbf{G}\textrm{O}(\infty,\infty)$ (respectively,
$\underset{\to}\lim GS(k'_m,V_{n'_m})=\mathbf{G}\textrm{S}(\infty,\infty)$). Without loss of generality we assume
that all $n_m$ in (\ref{exh1}) are odd.

Consider some $n_m$ such that $\mathbb{P}^q\subset GO(k_m,V_{n_m})$ (respectively, $\mathbb{P}^q\subset
GS(k'_m,V_{n'_m})$) and choose a maximal projective space $\mathbb{P}^r$ on $GO(k_m,V_{n_m})$ (respectively,
$GS(k'_m,V_{n'_m}$) such that $\mathbb{P}^q\subset\mathbb{P}^r$. The projective space $\mathbb{P}^r$ is either of type
$\mathbb{P}^r_{\alpha}$ or $\mathbb{P}^r_{\beta}$,  and we label $\mathbb{P}^q$ according to the label of
$\mathbb{P}^r$.
Lemma \ref{two possib},(i),(iii) (respectively, Lemma \ref{intersect dim=1},(i),(ii)) implies that this labeling is well
defined as long as the chain (\ref{exh1}) (respectively,  (\ref{exh2})) is fixed. Moreover, using Theorem 1 and
Lemma \ref{non-isom2} one can verify that the labelings $\mathbb{P}^q_{\alpha}$ and $\mathbb{P}^q_{\beta}$ are intrinsic
to the ind-variety $\mathbf{G}\mathrm{O}(\infty,\infty)$ (respectively, $\mathbf{G}\mathrm{S}(\infty,\infty)$), i.e. do
not depend on the
choice of chain (\ref{exh1}) (respectively,  (\ref{exh2})) satisfying the above conditions.

Let now $\mathbf{P}^{\infty}\hookrightarrow\mathbf{G}\mathrm{O}(\infty,\infty)$ (respectively,
$\mathbf{P}^{\infty}\hookrightarrow\mathbf{G}\mathrm{S}(\infty,\infty)$) be a linear embedding. We call its image an
\textit{infinite projective space $\mathbf{P}^{\infty}$ on} $\mathbf{G}\mathrm{O}(\infty,\infty)$ (respectively,
$\mathbf{G}\mathrm{S}(\infty,\infty)$). We say that $\mathbf{P}^{\infty}=\mathbf{P}^{\infty}_{\alpha}$ if
$\mathbf{P}^{\infty}=\underset{\to}\lim\mathbb{P}^q_{\alpha}$ for some projective spaces $\mathbb{P}^q_{\alpha}$ on
$\mathbf{G}\mathrm{O}(\infty,\infty)$ (respectively, $\mathbf{G}\mathrm{S}(\infty,\infty)$). In a similar way we define
$\mathbf{P}^{\infty}_{\beta}$ on $\mathbf{G}\mathrm{O}(\infty,\infty)$ (respectively,
$\mathbf{G}\mathrm{S}(\infty,\infty)$).

Next, we observe that Lemma \ref{point, line} implies that on $\mathbf{G}\mathrm{O}(\infty,\infty)$ there are pairs of
maximal infinite projective spaces $\mathbf{P}^{\infty}_{\alpha}$ and $\mathbf{P}^{\infty}_{\beta}$ such that
$\mathbf{P}^{\infty}_{\alpha}\cap\mathbf{P}^{\infty}_{\beta}$ is a point.

To complete the proof, we observe that on $\mathbf{G}\mathrm{S}(\infty,\infty)$ any two maximal infinite projective
spaces $\mathbf{P}^{\infty}_{\alpha}$ and $\mathbf{P}^{\infty}_{\beta}$ intersect in a projective line whenever their
intersection is non-empty. This follows from Lemma \ref{intersect dim=1}. More precisely, an infinite projective space
$\mathbf{P}^{\infty}_{\alpha}$ (respectively, $\mathbf{P}^{\infty}_{\beta}$) is maximal on $\mathbf{G}S(\infty,\infty)$
if and only if, for any chain (\ref{exh2}) the intersections
$\mathbf{P}^{\infty}_{\alpha}\cap GS(k'_m,V_{n'_mj})$ are maximal projective spaces in $GS(k'_m,V_{n'_m})$ for large
enough $m$. This is a consequence of Lemma \ref{two possib},(i). Now Lemma \ref{two possib},(iii) implies the assertion
that maximal projective spaces $\mathbf{P}^{\infty}_{\alpha}$ and $\mathbf{P}^{\infty}_{\beta}$ intersect in a
projective line whenever their intersection is non-empty.

Since the intersection properties of maximal infinite projective spaces $\mathbf{P}^{\infty}_{\alpha}$ and
$\mathbf{P}^{\infty}_{\beta}$ on $\mathbf{G}\mathrm{O}(\infty,\infty)$ and $\mathbf{G}\mathrm{S}(\infty,\infty)$
are intrinsic to the geometry of $\mathbf{G}\mathrm{O}(\infty,\infty)$ and $\mathbf{G}\mathrm{S}(\infty,\infty)$, we
conclude that $\mathbf{G}\mathrm{O}(\infty,\infty)$ and $\mathbf{G}\mathrm{S}(\infty,\infty)$ are non-isomorphic
ind-varieties.

The arguments in all other cases are similar. One either shows that on one of the ind-varieties in question there are
maximal projective spaces which do not exist on the other, or shows that the intersection properties of maximal
projective spaces are different on both ind-varieties. For instance, on $\mathbf{G}\mathrm{O}(k,\infty)$ there are
maximal projective spaces $\mathbb{P}^k_{\alpha}$ and $\mathbf{P}^{\infty}_{\beta}$ which intersect in a point, while on
$\mathbf{G}\mathrm{S}(k,\infty)$ two maximal projective spaces  $\mathbb{P}^k_{\alpha}$ and
$\mathbf{P}^{\infty}_{\beta}$ intersect in a projective line or do not intersect at all. We leave the details to the
reader.

\end{proof}



\vspace{1cm}





\begin{thebibliography}{MMMMM}

\bibitem[BV]{BV} W. Barth, A. Van de Ven, {\it On the geometry in codimension 2 in Grassmann
manifolds,} In: Lecture Notes in Math. 412, Springer-Verlag 1974, pp. 1-35.

\bibitem[DP]{DP} J. Donin, I. Penkov, {\it Finite rank vector bundles on inductive limits of
grassmannians}, IMRN No. 34 (2003), 1871-1887.

\bibitem[DiP]{DiP} I. Dimitrov, I. Penkov, {\it Ind-varieties of generalized flags as homogeneous
spaces for classical ind-groups}, IMRN  No. 55 (2004), 2935-2953.

\bibitem [H1]{H1} R. Hartshorne, {\it Stable reflexive sheaves}, Math. Ann. {\bf 254} (1980), 121-176.

\bibitem[OSS]{OSS}{ C. Okonek, M. Schneider, H. Spindler,} {\it Vector Bundles on Complex
Projective Spaces,} Birkh\"auser, 1980.

\bibitem[P]{P} I. Penkov, The Penrose transform on general Grassmannians, C. R. Acad. Bulg. des Sci. \textbf{22} (1980),
1439-1442.

\bibitem[PT1]{PT1} I. Penkov, A. S. Tikhomirov, {\it Rank-2 vector bundles on ind-Grassmannians,} Algebra, arithmetic, and geometry: in honor
of Yu. I. Manin, vol. II, Progr. Math., vol. 270, Birkhaeuser, Boston-Basel-Berlin 2009, pp. 555-572.

\bibitem[PT2]{PT2} I. Penkov, A. S. Tikhomirov, {\it Triviality of vector bundles on twisted ind-Grassmannians}.
Matematicheskij Sbornik, {\bf 202}, No.1 (2011), 65-104 (Russian). English translation:
Sbornik: Mathematics, {\bf 202}, No.1 (2011), 61-99.

\bibitem[PT3]{PT3} I. Penkov, A. S. Tikhomirov, {\it On the Barth-Van de Ven-Tyurin-Sato Theorem}, in progress.


\bibitem[S1]{S1} E. Sato, {\it On the decomposability of infinitely extendable vector bundles on
projective spaces and Grassmann varieties},  J. Math. Kyoto Univ. $\mathbf{17}$
(1977), 127-150.

\bibitem[S2]{S2} E. Sato, {\it On infinitely extendable vector bundles on $G/P$}, J. Math. Kyoto
Univ. $\mathbf{19}$ (1979), 171-189.

\bibitem[S3]{S3} E. Sato, {\it The decomposability of an infinitely extendable vector bundle on
the projective space},  Intern. Symp. on Algebraic Geometry, Kyoto, 1977, 663-672.

\bibitem[T]{T} A. N. Tyurin, {\it Vector bundles of finite rank over infinite varieties,} Math.
USSR Izvestija $\mathbf{10}$ (1976), 1187-1204.

\end{thebibliography}




\end{document}
